% =====================================================================
% IRH v93.1 — REBUILT TeX (NOT the author's source)
% Rebuilt 2026-11-04 from the PDF text of IRH v93.1
% (~/workspace/user/files/IRH_v93_1_1_44kz.pdf, 31 pp., and the
% extracted text ~/workspace/research/truth-seeking-v5/irh-test/irh_v93_1.txt).
% No .tex source for v93.1 was found in the workspace or on Google Drive
% (Drive search for IRH / .tex returned only PDFs, zips and an audit doc).
% This file is therefore a REBUILD, not a reissue of the author's source.
% Figures of the original are NOT recovered: figure environments below
% carry the original captions and are marked [figure placeholder].
% "Expansion (rebuild)" passages were added in this rebuild: they show
% step-by-step arithmetic that the source text skips. They add no new
% physical claim. Corrections applied in this rebuild are listed in the
% front matter. Nothing in this rebuild upgrades any claim's HLRE grade.
% =====================================================================
\documentclass[11pt,a4paper]{article}
\usepackage{geometry}
\geometry{margin=2.5cm}
\usepackage{fontspec}
\usepackage{unicode-math}
\usepackage{amsmath}
\usepackage{booktabs}
\usepackage{xcolor}
\usepackage{hyperref}
\hypersetup{colorlinks=true,linkcolor=blue!60!black,urlcolor=blue!60!black}

% ---- FONT RULE (exact, per rebuild brief) ----
% Fira Sans for text, Fira Mono for code, Fira Math for math (Scale=1.15).
% No fake-bold emulation is used anywhere in this file.
\setmainfont{Fira Sans}
\setsansfont{Fira Sans}
\setmonofont{Fira Mono}
\setmathfont{Fira Math}[Scale=1.15]
% Fallback for missing glyphs ONLY (cmap evidence, 2026-11-04, fontTools):
%  U+2208 (in)      : PRESENT in Fira Math -> needs no fallback; math \in used.
%  U+2216 (setminus): ABSENT in Fira Math, PRESENT in Noto Sans Math -> fallback:
\setmathfont{Noto Sans Math}[range={"2216},Scale=1.15]
%  U+2090 (subscript a): ABSENT in Fira Sans, Fira Math and Noto Sans Math,
%  PRESENT in Noto Sans -> text-mode fallback via an explicit font switch
%  ({\notosans <glyph>}) at the single site where the glyph is used:
\newfontfamily{\notosans}{Noto Sans}
% Text-mode fallback note: a literal U+2208 in text would also fall back via
% math mode; body text uses $\in$ and $\setminus$ in math mode throughout.

\newcommand{\expansion}[1]{\par\smallskip\noindent\textbf{Expansion (rebuild).} #1\par\smallskip}
\newcommand{\flagbox}[1]{\par\smallskip\noindent\textbf{Flag (rebuild).} #1\par\smallskip}
\newcommand{\gradebox}[1]{\par\smallskip\noindent\textbf{Grade note.} #1\par\smallskip}

\title{Intrinsic Resonance Holography:\\Dynamical Selection of a $D_4$ Substrate from Interference Dynamics\\[0.4em]\large Version 93.1 (rebuilt \TeX)}
\author{Brandon D. McCrary\\Independent Theoretical Physics Researcher\\ORCID: 0009-0008-2804-7165}
\date{September 17, 2026 (original); rebuilt November 4, 2026}

\begin{document}
\maketitle

% Font-fallback demonstration line (rebuild apparatus, not manuscript content):
\noindent\textit{Font check (rebuild apparatus):} $x \in A \setminus B$, subscript-a glyph {\notosans ₐ}, $D_4$, $\tau$, $q_4$. All three required glyphs (U+2208, U+2216, U+2090) render in this document: $\in$ from Fira Math, $\setminus$ from the Noto Sans Math fallback, {\notosans ₐ} from the Noto Sans text fallback.

\section*{Rebuild front matter --- read first}
\addcontentsline{toc}{section}{Rebuild front matter}

\textbf{Source status: REBUILT FROM PDF, not the author's source.} No \texttt{.tex} file for IRH v93.1 exists in the workspace ($\sim$/workspace/research/irh/tex/ holds only arXiv reference downloads) or on the connected Google Drive (searched 2026-11-04: Drive queries for \texttt{IRH} and for \texttt{.tex} returned only \texttt{IRH\_v93\_1.pdf}, verification zips, and an audit document). The manuscript below was reconstructed from the PDF text of v93.1 (31 pages, created with \LaTeX{} and hyperref, September 18, 2026). Wording follows the extracted text closely; where the PDF text extraction garbled a formula, the mathematically verified form is used and the garble is noted.

\textbf{Corrections applied in this rebuild (each flagged again at the site):}
\begin{enumerate}
\item \textbf{$\tau_{D_4}/\tau_{\mathrm{BCC}} = \sqrt{2}$, NOT $2$.} The source states the ratio as $2$ (abstract, \S II.5, \S VI-A.1). That is arithmetically wrong: $\tau_{D_4} = 4/\sqrt{6}$ and $\tau_{\mathrm{BCC}} = 2/\sqrt{3}$, so the ratio is $(4/\sqrt{6})/(2/\sqrt{3}) = \sqrt{2}$. Verified in SymPy 1.14.0 and in python-flint 0.8.0 (squared ratio $= 2$). Every occurrence in this rebuild reads $\sqrt{2}$.
\item \textbf{Ordered-branch $r_c$ form.} The source's \S II.3 gives $B_* = 2u\tau/(3vq_4)$ and $r_c = 2u^2\tau^2/(9vq_4)$; the claim ledger (\S VIII.1) as extracted reads ``$r_c = 92\,u^2\tau^2/(vq_4)$'', which is a PDF-text garble of $2/9$. This rebuild uses the \S II.3 form, verified symbolically in SymPy.
\item \textbf{Phase-band percentage inconsistency flagged, not resolved.} \S II.7 states the band occupies 11.85\% of the scanned plane, then states the $w_6 = 0$ value as 11.90\%. The 0.05-point discrepancy is preserved and flagged at the site; neither value was silently overwritten.
\item \textbf{Dimension wording clarified at \S IV.1.} The source's ``dimension increases $15 \to 21$'' refers to $\dim(\mathfrak{g}_2 \oplus \mathfrak{u}(1)) = 14 + 1 = 15$ versus $\dim(\text{Pati--Salam}) = 15 + 3 + 3 = 21$, not to $\dim \mathfrak{so}(8) = 28$. This rebuild states both explicitly.
\end{enumerate}

\textbf{Grade discipline (HLRE v5; this rebuild changes no grade).} The $D_4$/triality/$2T$--McKay--$E_6$/Swift--Hohenberg \emph{mathematics} is Grade~A. IRH's \emph{physical/substrate reading} (``$D_4$ is the physical substrate'', the gauge reading as physics) is graded $\leq$ C/D. Three generations are \textbf{not derived}: they are an acknowledged INPUT (\S V.2). A SymPy check of arithmetic in this rebuild is \textbf{not} a grade upgrade for any physical interpretation.

\textbf{Two quantities that must never be identified.} (i) Perelman's $\tau$ (backward time/temperature in Ricci-flow theory) is \textbf{not} IRH's $\tau$ (here $\tau \equiv T/N^{3/2}$, counting ordered zero-sum triads on the reciprocal shell). (ii) IRH's $S = \tau^2/q_4 = 64/141$ is \textbf{shell combinatorics}, not an entropy, and must not be identified with Perelman's $W$-functional. Both distinctions are repeated at the sites where confusion is possible.

\textbf{What was verified, and what was flagged unverifiable-from-source.} Every numerical/algebraic claim checkable from the source text alone was verified with SymPy 1.14.0 in the conda environment \texttt{jacobian-conjecture-2d}; load-bearing claims were additionally checked with a second independent tool (python-flint 0.8.0, Singular 4.4.1, or exhaustive Python enumeration). The full log is \texttt{sympy\_verify.log} (59 PASS, 1 intentional FAIL-flag: the 11.85\%/11.90\% inconsistency), \texttt{second\_tool\_verify.log} (20 PASS), \texttt{singular\_verify.log} (moment identities $= 0$, $\det E_6 = 3$, $64/141$), and \texttt{sympy\_supplement.log} (7 PASS). Claims that rest on simulations, scans, or enumerations whose code and data are not in the source text are marked \textbf{FLAG: simulation/scan claim, not re-run in this leg} rather than given invented numbers. A complete flagged list is in the rebuild addendum (the rebuild's own closing section, after Part X of the manuscript).

\begin{abstract}
Intrinsic Resonance Holography proposes that the structure of physical reality rests on a self-organizing oscillatory substrate whose stable interference pattern is the $D_4$ root lattice. This document is the complete consolidated statement of the framework following five independent verification cycles, in which every load-bearing claim was either re-derived by executed computation, validated against a known limiting case, confirmed by simulation, excluded by theorem, or --- new in this version --- machine-checked in Lean~4.

The framework in this version differs from v92 in kind, not merely in detail. Its central ontological commitment has changed: $D_4$ is no longer a postulated static lattice but the nodal interference pattern of a field obeying tension--relaxation dynamics. That reformulation removes a primitive, dissolves a standing circularity objection at its root rather than patching it, and supplies the selection mechanism the static formulation never possessed.

What is now established. The $D_4$ pattern is selected dynamically rather than assumed, by a parameter-free criterion $S = \tau^2/q_4$ derived symbolically, validated by recovering the Alexander--McTague BCC result in three dimensions~\cite{alexander1978}, applied in four where $D_4$ wins at $S = 0.454$ with selection quantities that are exact rationals and simple radicals ($64/141$, $16/45$, $73/423$, $\tau_{D_4}/\tau_{\mathrm{BCC}} = \sqrt{2}$ \textbf{[corrected from $2$; see front matter]}) rather than fitted decimals, confirmed by direct simulation from random initial conditions, and independently corroborated by a published molecular-dynamics result obtained by an unrelated method~\cite{hoy2024}. An independent mechanism-level reproduction finds 77.1\% of spectral power on $D_4$ with triads active versus 21.1\% with triads removed, rising to 89.5\% under annealing --- establishing the cubic triad term as the causal selection mechanism. The full 24-mode standing wave is a genuine dynamical attractor with relaxed free energy $F/V = -0.0077276$, reproduced independently at $-0.007728$. The resonance maximum is unique modulo lattice periodicity (32 periodic maxima on the $16^4$ grid); the transition is first order. The lattice constant $a_0 = 2\pi/q_0$ is derived rather than postulated. Four-dimensionality acquires a new structural argument: the Swift--Hohenberg~\cite{swift1977} curvature term is exactly scale-invariant at $d = 4$. The gauge algebra is identified as $E_6$, read off the pattern's own group structure via the binary tetrahedral group and the McKay correspondence~\cite{mckay1980} rather than imposed to fit a target. The geometric core is now machine-checked: eight Lean~4 modules compile cleanly with no sorry, admit, or axiom.

What is excluded by theorem. $SO(8)$ cannot be the gauge algebra, by two independent proofs. The Standard Model algebra does not embed in $\mathfrak{so}(8)$ by any route. The anisotropy coefficient $\xi_2 = -1/384$ is not derivable from the stated lattice. Both time-step values stated in v92 violate the framework's own stability bound. Three generations is not a derived result: five candidate mechanisms were tested and all five are closed. Monopole absence, claimed in earlier versions, is inconsistent with an $E_6$ gauge group.

What remains open is stated exhaustively in \S IX, classified by the type of work each item requires. Every correction made during verification --- including corrections to the author's own prior output --- is recorded in the appendices, and every externally attributable claim carries a linked numerical citation.

The framework claims substantially less than v92 did. What it now claims has been checked --- by hand, by simulation, by independent reproduction, and now by machine.
\end{abstract}

\section*{How to Read This Document}
Part I states the ontological reformulation and its immediate consequences. Parts II--III present the geometric and dynamical core, which contains the strongest results. Part IV presents the algebraic sector, including the exclusion theorem and the gauge decision. Part V covers matter and defects, including the negative result on generations. Part VI lists falsifiable predictions. Part VII states retractions from v92. Parts VIII--IX give the complete claim ledger and open-problem register. Appendices record every correction made during verification, including corrections to this author's own prior output.

Epistemic labels used throughout:
\begin{itemize}
\item \textbf{DERIVATION} --- follows from stated assumptions with no empirical input
\item \textbf{PREDICTION} --- falsifiable model output, independent of the data compared against
\item \textbf{CALIBRATION} --- a parameter adjusted to match data; never presented as a prediction
\item \textbf{INPUT} --- an assumption, explicitly acknowledged as such
\item \textbf{EXCLUDED} --- ruled out by proof
\item \textbf{OPEN} --- a well-posed question not resolved here
\end{itemize}
A claim carrying no label has not been graded and should be treated as unsupported.

\gradebox{IRH's $\tau$ is defined in \S II.3 as $T/N^{3/2}$ over ordered zero-sum triads on the reciprocal shell. It is not Perelman's $\tau$, and $S = \tau^2/q_4$ is not Perelman's $W$. See the front matter.}

% ============================ PART I ============================
\section{PART I --- THE ONTOLOGICAL REFORMULATION}

\subsection{I.1 Why the Static Formulation Failed}
Version 92 treated $D_4$ as a static four-dimensional lattice with fields defined at sites, $\phi(n, s)$ for $n \in \Lambda_{D_4}$. Three independent objections were raised against that formulation and all three are sustained.

\begin{figure}[h]
\centering
\fbox{\parbox{0.85\textwidth}{\centering [figure placeholder --- original Figure 1 not recovered from PDF text]\\The logical architecture of the framework. Arrows indicate dependency: an item cannot be established before those above it. The forcing result at the centre is what compelled the introduction of internal structure and therefore the gauge decision of Part IV.}}
\caption{Figure 1. The logical architecture of the framework.}
\end{figure}

\textbf{The primitives were circular.} The framework designated $a_0$ and $\Omega_P$ as primitive but defined them as $a_0 = L_P/\sqrt{24}$ and $\Omega_P = \sqrt{24}\,c/L_P$ with $L_P = \sqrt{\hbar G/c^3}$. These are re-expressions of measured Planck-scale quantities, not independent microscopic inputs. The triad relations $c = a_0\Omega_P$, $\hbar = M_* c a_0$, $G = 24c^2a_0/M_*$ are algebraically consistent once the definitions are chosen, but a change of variables is not a derivation of the constants.

\textbf{The lattice was not selected.} No mechanism explained why the substrate would adopt $D_4$ periodicity rather than $\mathbb{Z}^4$ or any other lattice. $D_4$ was chosen because it fits, which is reverse engineering, not derivation.

\textbf{The action was ill-defined.} The formulation mixed index-lattice and dimensionful coordinates, writing shifts such as $\phi(n + a_0 e_\mu)$ for $n \in \mathbb{Z}^4$, and stated two mutually inconsistent time steps differing by $2\pi$.

\subsection{I.2 The Reformulation}
\textbf{Statement.} The primitive is an oscillatory tension field $\psi(x, t)$ on $\mathbb{R}^4$. $D_4$ is not postulated; it is the nodal interference pattern into which the field self-organizes. Persistence is dynamic rather than static: the pattern is a standing wave maintained by a balance of tension and relaxation. In the framework's own formulation, the primitive is not a noun but a verb.

\textbf{The dynamical equation.} Tension and relaxation together specify a damped wave equation:
\begin{equation}
\partial_t^2 \psi + \gamma \partial_t \psi = -\frac{\delta F}{\delta \psi},
\qquad
F[\psi] = \int d^4x \left[ \tfrac{1}{2}\psi\left(r + (q_0^2 + \nabla^2)^2\right)\psi - \tfrac{u}{3}\psi^3 + \tfrac{v}{4}\psi^4 \right].
\end{equation}
Here $\partial_t^2\psi$ is the tension term, $\gamma\partial_t\psi$ the relaxation term, and $(q_0^2 + \nabla^2)^2$ the Swift--Hohenberg operator~\cite{swift1977} penalizing wavevectors away from the resonant shell $|k| = q_0$. The parameter $r$ is the quench depth and $u, v$ the nonlinear couplings.

\subsection{I.3 Immediate Consequences}
\textbf{DERIVATION --- the lattice constant is no longer primitive.}
\begin{equation}
a_0 = \frac{2\pi}{q_0}.
\end{equation}
The lattice constant is fixed by the substrate's resonant wavenumber. This is not a restatement of the old circularity in new symbols: $a_0$ is now an output of the substrate's dispersion rather than an input defined through $L_P$. The objection of \S I.1 dissolves rather than being patched.

\textbf{DERIVATION --- the site-field action is retired.} A field defined at sites presupposes the static ontology. Under the reformulation, sites are not carriers of degrees of freedom; they are loci where $|\psi|$ is extremal. The field is $\psi$ and $D_4$ is a property of its solution. Results derived from the site-field action (the CFL bound, the discrete dispersion relation) remain valid for that formulation and are reported in Part III with that scope stated.

\textbf{DERIVATION --- the substrate must carry conserved dynamics.} Verified numerically: under non-conserved (Model A~\cite{hohenberg1977}) dynamics $\partial_t\psi = -\delta F/\delta\psi$, the mean field drifts by $1.8 \times 10^{-1}$ over a run; under conserved (Model B~\cite{hohenberg1977}) dynamics $\partial_t\psi = \nabla^2(\delta F/\delta\psi)$ it is preserved to $8.3 \times 10^{-17}$. Since Model A drives the baseline tension out of the band in which the full $D_4$ pattern is stable (\S II.7), conserved dynamics is a requirement of the framework, not a modeling choice. The substrate's total tension must be a conserved quantity.
\flagbox{The drift figures $1.8\times10^{-1}$ and $8.3\times10^{-17}$ are reported simulation results. The simulation code is not in the source text, so these numbers are \textbf{not re-run in this leg}; they are carried as stated, not as re-verified.}

% ============================ PART II ============================
\section{PART II --- THE GEOMETRIC CORE}

\subsection{II.1 Mathematical Properties of $D_4$}
\textbf{DERIVATION --- exact isotropy through fourth order.} Enumerating all 24 roots $\Phi_{D_4} = \{\pm e_i \pm e_j\}$ and summing symbolically:
\begin{equation}
\sum_{\alpha \in \Phi_{D_4}} (k \cdot \alpha)^2 = 12|k|^2,
\qquad
\sum_{\alpha \in \Phi_{D_4}} (k \cdot \alpha)^4 = 12|k|^4.
\end{equation}
Both hold exactly, with zero symbolic residual and zero directional dependence, confirmed numerically across 400 random directions at a relative spread of $3 \times 10^{-16}$. This is guaranteed rather than coincidental: the $D_4$ root system is a spherical 5-design~\cite{delsarte1977}, so any polynomial of degree $\leq 5$ averaged over the 24 points equals its average over the continuous sphere, and $(k\cdot\hat{n})^4$ has degree 4. Genuine anisotropy first appears at sixth order, consistent with the fact that no 24-point configuration can form a spherical 6-design~\cite{delsarte1977}.

\expansion{Per pair $(i,j)$ the four roots contribute $(k_i+k_j)^2 + (k_i-k_j)^2 + (-k_i+k_j)^2 + (-k_i-k_j)^2 = 4k_i^2 + 4k_j^2$ at second order, and $4k_i^4 + 24k_i^2k_j^2 + 4k_j^4$ at fourth order. Each index $i$ occurs in 3 pairs, so the totals are $12\sum_i k_i^2 = 12|k|^2$ and $12\sum_i k_i^4 + 24\sum_{i<j} k_i^2k_j^2 = 12|k|^4$. SymPy 1.14.0 confirms both identities symbolically with zero residual; Singular 4.4.1 confirms both as exact polynomial identities (differences $0$). The 400-direction numerical spread $3\times10^{-16}$ is a reported result, \textbf{not re-run in this leg}.}

\begin{figure}[h]
\centering
\fbox{\parbox{0.85\textwidth}{\centering [figure placeholder --- original Figure 2 not recovered]\\The fourth moment evaluated at 400 random directions on $S^3$, for the true $D_4$ root system (24 terms) against the hypercubic stand-in (8 terms) from which the claimed anisotropy coefficient was derived. The $D_4$ sum is constant to machine precision; the hypercubic sum varies by 32\% with direction. These are not the same calculation.}}
\caption{Figure 2. Fourth moment: $D_4$ versus hypercubic stand-in.}
\end{figure}

\textbf{DERIVATION --- isoduality.} The reciprocal minimal shell of the $D_4$ lattice contains 24 vectors whose pairwise-angle fingerprint $\{-1, -\tfrac{1}{2}, 0, \tfrac{1}{2}\}$ is identical to that of the root system itself. $D_4$ is an isodual lattice~\cite{conway1999}, geometrically congruent to its own dual --- a property shared by only a short list of lattices.

\subsection{II.2 The Resonance-Cavity Mechanism}
The framework's informal cymatics picture is a well-posed reciprocal-lattice problem: which set of plane-wave modes, superposed, interferes constructively exactly at the points of $D_4$?

\textbf{DERIVATION.} Compactifying the substrate on the torus $\mathbb{R}^4/\Lambda$ with $\Lambda \cong D_4$ forces the allowed mode content to be the reciprocal shell. Isoduality then makes that forced content congruent to $D_4$ itself. Building $\Psi(x) = \sum_g \cos(g \cdot x)$ from the shell and evaluating exhaustively over $25^4 = 390{,}625$ points of one fundamental domain:
\begin{center}
\begin{tabular}{ll}
\toprule
test point & $\Psi$ \\
\midrule
a genuine $D_4$ lattice point & 24.0 (global maximum, unique) \\
nearest $\mathbb{Z}^4$ point not in $D_4$ & $-8.0$ \\
deep hole $(\tfrac12,\tfrac12,\tfrac12,\tfrac12)$ & $-8.0$ \\
\bottomrule
\end{tabular}
\end{center}
No other point came within 0.1\% of the maximum. Periodicity under $D_4$ translations was confirmed to machine precision.

\expansion{$25^4 = 390{,}625$ exactly. The count 24.0 equals the shell size (all 24 cosines equal 1 at a lattice point); $-8.0$ at the other two test points is the stated scan outcome.}
\flagbox{The exhaustive scan itself is a reported computation; it is \textbf{not re-run in this leg} (no scan code in the source text).}

\textbf{Scope of ``unique.''} The maximum is unique within one fundamental domain, which is where the scan was performed. On a full periodic grid it necessarily recurs at every lattice image: on a $16^4$ grid carrying two pattern periods per direction it is attained at 32 points.
\expansion{Two periods per direction give $2^4 = 16$ pattern repeats; the $D_4$ checkerboard contributes two sublattices (corner and body-centre families), so $16 \times 2 = 32$, not the naive 16. The claim should always be stated as unique modulo lattice periodicity.}
No external boundary is required: the periodicity is intrinsic to the medium, so the substrate is simultaneously medium and boundary.

\textbf{Scope limit.} Isoduality alone does not select $D_4$ over competitors --- $\mathbb{Z}^4$ is also isodual and self-resonates trivially. Selection requires the dynamics of \S II.3.

\subsection{II.3 The Selection Functional and Criterion}
\textbf{DERIVATION.} Expand the field on the minimal reciprocal shell with $N$ modes of equal amplitude, writing $T$ for the ordered triples summing to zero (closed wavevector triads) and $Q$ for the ordered quadruples. In the normalized amplitude $B = A\sqrt{N}$:
\begin{equation}
F(B) = \frac{r}{2}B^2 - \frac{u\tau}{3}B^3 + \frac{vq_4}{4}B^4,
\qquad
\tau \equiv \frac{T}{N^{3/2}}, \quad q_4 \equiv \frac{Q}{N^2}.
\end{equation}
A first-order transition occurs where the ordered minimum becomes degenerate with the disordered state, $F = 0$ and $F' = 0$ simultaneously. Solving that system symbolically:
\begin{equation}
B_* = \frac{2u\tau}{3vq_4}, \qquad r_c = \frac{2u^2\tau^2}{9vq_4}.
\end{equation}
The pattern ordering first --- at the highest $r$, meaning the shallowest quench --- maximizes
\begin{equation}
S = \frac{\tau^2}{q_4}.
\end{equation}
This criterion is parameter-free in $u$ and $v$: their values scale $r_c$ but cannot reorder the candidates. Selection depends only on the combinatorics of the reciprocal shell.

\expansion{Solving $F = 0$ and $F' = 0$ step by step (SymPy 1.14.0, symbolic in $B, r, u, v, \tau, q_4$): from $F' = 0$ with $B \neq 0$, $r - u\tau B + vq_4 B^2 = 0$; from $F = 0$ with $B \neq 0$, $\tfrac{r}{2} - \tfrac{u\tau}{3}B + \tfrac{vq_4}{4}B^2 = 0$. Multiply the second equation by 2: $r - \tfrac{2u\tau}{3}B + \tfrac{vq_4}{2}B^2 = 0$. Subtract from the first: $-\tfrac{u\tau}{3}B + \tfrac{vq_4}{2}B^2 = 0$, hence $B_* = \tfrac{2u\tau}{3vq_4}$. Substituting back: $r_c = u\tau B_* - vq_4 B_*^2 = \tfrac{2u^2\tau^2}{3vq_4} - \tfrac{4u^2\tau^2}{9vq_4} = \tfrac{2u^2\tau^2}{9vq_4}$. Since $r_c \propto \tau^2/q_4 = S$ with the common factor $2u^2/(9v)$, the ordering by $r_c$ is exactly the ordering by $S$, independent of $u$ and $v$.}
\gradebox{$S$ here is shell combinatorics ($\tau$ counts triads, $q_4$ counts quadruples). It is not an entropy and not Perelman's $W$. IRH's $\tau$ is not Perelman's $\tau$.}

\subsection{II.4 Validation in Three Dimensions}
Before application in four dimensions the criterion was required to reproduce the known result.
\begin{center}
\begin{tabular}{llllll}
\toprule
crystal & $N$ & $T$ & $\tau$ & $q_4$ & $S = \tau^2/q_4$ \\
\midrule
BCC & 12 & 48 & 1.1547 & 3.750 & 0.35556 \\
FCC & 8 & 0 & 0 & 3.375 & 0 \\
SC & 6 & 0 & 0 & 2.500 & 0 \\
\bottomrule
\end{tabular}
\end{center}
BCC is selected; FCC and SC are excluded by vanishing triad content --- their reciprocal minimal shells contain no closed triangles. This is the Alexander--McTague result~\cite{alexander1978}, recovered from first principles by the same code subsequently applied in four dimensions.

\expansion{BCC step by step: $\tau = 48/12^{3/2} = 48/(12 \cdot 2\sqrt{3}) = 2/\sqrt{3} = 1.1547005\ldots \approx 1.1547$; $\tau^2 = 4/3$; $q_4 = 15/4 = 3.750$ (so $Q = q_4N^2 = 540$); $S = (4/3)/(15/4) = 16/45 = 0.35555\ldots \approx 0.35556$. FCC: $q_4 = 27/8 = 3.375$ (so $Q = 216$), $T = 0 \Rightarrow S = 0$. SC: $q_4 = 5/2 = 2.500$ (so $Q = 90$), $T = 0 \Rightarrow S = 0$. SymPy and python-flint agree on all values.}
\textbf{Scope nuance.} The BCC preference is a near-melting result; Groh and Mulder~\cite{groh1999} showed that symmetry breaking at the liquid--solid transition need not produce bcc outside that regime. The $S$ criterion reproduces the Alexander--McTague conclusion within its domain of validity.

\subsection{II.5 Result in Four Dimensions}
\textbf{PREDICTION.}
\begin{center}
\begin{tabular}{lllllll}
\toprule
lattice & $N$ & $T$ & $\tau^2$ (exact) & $q_4$ (exact) & $S = \tau^2/q_4$ (exact) & decimal \\
\midrule
$D_4$ & 24 & 192 & $8/3$ & $47/8$ & $64/141$ & 0.453901 \\
$A_4^*$ & 20 & 120 & $9/5$ & $93/20$ & $12/31$ & 0.387097 \\
$\mathbb{Z}^4$ & 8 & 0 & 0 & $21/8$ & 0 & 0 \\
$A_4$ & 10 & 0 & 0 & $27/10$ & 0 & 0 \\
\bottomrule
\end{tabular}
\end{center}
\textbf{Table convention.} Rows are labelled by the crystal; $N$ is the size of that crystal's reciprocal minimal shell, since the density is expanded in reciprocal vectors. Hence the $A_4^*$ row shows $N = 20$ (the $A_4$ shell) and the $A_4$ row shows $N = 10$ (the $A_4^*$ shell). Both are correct under this convention, which was not stated in earlier versions and caused a reasonable query in independent review.

The selection quantities are exact rationals, not fitted decimals: $S_{D_4}/S_{A_4^*} = 496/423$, a margin of $73/423 \approx 17.26\%$; $S_{\mathrm{BCC}} = 16/45$; and $\tau_{D_4}/\tau_{\mathrm{BCC}} = \sqrt{2}$ exactly \textbf{[corrected; the source reads $2$]}. $D_4$ is selected by a 17.3\% margin over the nearest genuine competitor. The hypercubic lattice and $A_4$ are structurally excluded: with $T = 0$ they cannot order through the cubic mechanism at all, exactly as FCC cannot in three dimensions.

Closed forms: $\tau_{D_4} = 4/\sqrt{6}$ and $\tau_{\mathrm{BCC}} = 2/\sqrt{3}$, in ratio exactly $\sqrt{2}$ \textbf{[corrected]}. $D_4$ is the four-dimensional analogue of BCC, and a stronger one.

\expansion{$D_4$ step by step: $\tau = 192/24^{3/2} = 192/(24\cdot 2\sqrt{6}) = 4/\sqrt{6}$; $\tau^2 = 16/6 = 8/3$; $q_4 = Q/N^2 = 3384/576 = 47/8$ (divide by 72: $3384/72 = 47$, $576/72 = 8$); $S = (8/3)/(47/8) = 64/141 = 0.4539007\ldots \approx 0.453901$. $A_4^*$: $\tau = 120/20^{3/2} = 3/\sqrt{5}$, $\tau^2 = 9/5$; $q_4 = 93/20$ (so $Q = 1860$); $S = (9/5)/(93/20) = 36/93 = 12/31 = 0.3870967\ldots \approx 0.387097$. Ratio: $(64/141)/(12/31) = (64\cdot 31)/(141\cdot 12) = 1984/1692 = 496/423$ (divide by 4); margin $496/423 - 1 = 73/423 = 0.1725768\ldots \approx 17.26\%$. $\tau$ ratio: $(4/\sqrt{6})/(2/\sqrt{3}) = 2\sqrt{3}/\sqrt{6} = 2/\sqrt{2} = \sqrt{2}$; squared ratio $(8/3)/(4/3) = 2$, confirming $\sqrt{2}$ and refuting $2$. SymPy 1.14.0, python-flint 0.8.0, and Singular 4.4.1 (rational check) all agree. $\mathbb{Z}^4$: $Q = (21/8)\cdot 64 = 168$, $T = 0 \Rightarrow S = 0$. $A_4$: $Q = (27/10)\cdot 100 = 270$, $T = 0 \Rightarrow S = 0$. These $Q$ recoveries check the table's internal arithmetic; the geometric enumeration of $Q$ for $A_4^*$, $\mathbb{Z}^4$, and $A_4$ is stated in the source, not re-derived here, and is flagged as such.}
\flagbox{The $Q$ values 1860, 168, 270 are consistency recoveries ($Q = q_4N^2$), not independent shell enumerations. The source's $D_4$ value $Q = 3384$ is used in \S II.8 and is consistent with $q_4 = 47/8$; its own enumeration is stated, not re-run.}

\textbf{Independent corroboration.} This analytic prediction agrees with a published molecular-dynamics study in which supercooled four-dimensional Lennard-Jones liquids were observed to crystallize into $D_4$ rather than the competing $A_4$ lattice~\cite{hoy2024}. At zero pressure and $0.59 < T < 0.63$, the liquid crystallizes within $\sim 2 \times 10^4\tau$; with 8 atoms per $D_4$ unit cell, the choice $N = 648$ promotes a $3 \times 3 \times 3 \times 3$-cell $D_4$ single crystal. The two determinations share no method: one is a Landau~\cite{landau1937} free-energy argument over reciprocal-shell combinatorics, the other a direct particle simulation under a standard pairwise potential. Compatible freezing behavior is also seen in four-dimensional hard spheres, where the densest Voronoi cluster is consistent with $D_4$ symmetry~\cite{vanmeel2009}.

\textbf{Prior art in four dimensions.} The Alexander--McTague Landau argument has been generalized to higher dimensions: Elitzur, Karni, and Rabinovici~\cite{elitzur2011} showed that in 4D the cross-polytope (16-cell) maximizes the cubic term, with the reciprocal lattice generated by $(\pm1,\pm1,\pm1,\pm1)$ --- the $D_4$ lattice. Theirs is an equilibrium near-melting argument; the $S$ criterion here is dynamical, selecting from noise by interference. The two routes are independent and select the same winner.

\textbf{Structural unification.} The 192 triads arise from exactly the 24 reciprocal vectors of \S II.2 --- same fingerprint, 8 partners at $120^\circ$ from each vector, $24 \times 8 = 192$. The interference picture and the energy picture are the same computation viewed twice. They interlock rather than merely agree.
\expansion{The ``8 partners'' claim was checked by exhaustive enumeration in this rebuild: for each of the 24 roots $\pm e_i \pm e_j$ (norm$^2 = 2$), exactly 8 other roots have dot product $-1$, i.e.\ $\cos\theta = -1/2$, $\theta = 120^\circ$. Hence $24 \times 8 = 192$ ordered triad partners, matching $T = 192$.}

\subsection{II.6 Dynamical Confirmation by Simulation}
Minimizing a functional and reaching its minimum dynamically are different claims.

\textbf{Method and full parameters.} The damped-wave equation was integrated on a $16^4$ periodic grid from random initial conditions using a semi-implicit pseudo-spectral scheme~\cite{trefethen2000}, unconditionally stable with the linear operator treated implicitly in Fourier space.

\textbf{Domain and normalization.} Wavevectors are integer-valued, $k = n$, with the box $L = 2\pi$ so that physical $k = 2\pi n/L = n$. The $D_4$ shell is the scaled root system $\{(\pm2,\pm2,0,0)$ and perms$\}$, giving $|k|^2 = 8$; this places two pattern periods per direction in the box. A reader using the unscaled root normalization $|\alpha|^2 = 2$ rescales accordingly; the two are the same physical configuration.

\textbf{The competitor shell, specified.} With $q_0^2 = 7.5$, the two competing shells are
\begin{center}
\begin{tabular}{lllll}
\toprule
shell & $|k|^2$ & modes & $T$ & linear penalty $(q_0^2 - |k|^2)^2$ \\
\midrule
$D_4$ star & 8 & 24 & 192 & 0.25 \\
competitor & 7 & 64 & 0 & 0.25 \\
\bottomrule
\end{tabular}
\end{center}
The competitor is the full integer shell at $|k|^2 = 7$. It carries zero triad content and 2.7 times as many modes, and sits at an identical linear penalty. The linear growth rates are therefore exactly degenerate and only the cubic term can break the tie. Remaining parameters: $r = -0.35$, $u = v = 1$, $\gamma = 1.2$, $\Delta t = 0.05$. 6000 steps direct quench; annealed run ramps $r$ from $-0.05$ to $-0.45$ over 14000 steps with decaying stochastic forcing.

\expansion{Shell counts: $D_4$ scaled shell --- choose the 2 nonzero positions in $\binom{4}{2} = 6$ ways, 4 sign choices, $6 \times 4 = 24$ modes, $|k|^2 = 4 + 4 = 8$. Competitor $|k|^2 = 7$ shell --- vectors $(\pm2,\pm1,\pm1,\pm1)$: 4 positions for the 2, $2^4 = 16$ sign choices, $4 \times 16 = 64$ modes. Penalties: $(7.5-8)^2 = 0.25$ and $(7.5-7)^2 = 0.25$ --- exactly degenerate, as claimed. Mode ratio $64/24 = 2.666\ldots \approx 2.7$.}

\textbf{Result (author's run).} The $D_4$ shell captured 84\% of spectral power on direct quench and 100\% after annealing. The control run with the cubic coupling set to zero produced a near-tie, 44\% against 56\%. The triad term is causally responsible for the selection.

\textbf{Independent reproduction.} A mechanism-level reproduction from an independent implementation --- different random seeds, random-field initialization of standard deviation 0.05, annealing noise amplitude 0.5 with linear decay --- found 77.1\% of spectral power on $D_4$ with triads active versus 21.1\% with triads removed, rising to 89.5\% under annealing and still increasing at step 14000. Exact percentages are implementation-dependent; the causal conclusion is identical: the cubic triad term reverses the selected winner.
\flagbox{All percentages in this section are reported simulation outcomes (author's run and independent reproduction). No simulation code is in the source text, so none was re-run in this leg. The causal reading (triads flip the winner) is carried as the source states it; the numbers themselves are not re-verified here.}

\begin{figure}[h]
\centering
\fbox{\parbox{0.85\textwidth}{\centering [figure placeholder --- original Figure 3 not recovered]\\Left: spectral power with the triad coupling on versus off. Right: the parameter-free ordering criterion in four dimensions, with the validated three-dimensional BCC value marked for scale.}}
\caption{Figure 3. Triad on/off spectral power; 4D ordering criterion.}
\end{figure}

\subsection{II.7 Baseline Tension, and the Phase Diagram}
The annealed run populated only 6 of 24 modes, prompting a comparison the functional analysis had not performed: against sub-patterns within the shell rather than only against other lattices.

\textbf{An initial failure.} At zero baseline field a 6-mode hexagonal sub-pattern has lower free energy than the full 24-mode pattern at every quench depth tested. The simulation was right and the analysis had been incomplete.

\textbf{Resolution.} The omitted variable is the substrate's baseline tension $\bar{\psi}$. Expanding $\psi = \bar{\psi} + A\sum_k e^{ik\cdot x}$ and collecting symbolically:
\begin{equation}
F_{\mathrm{mod}} = \frac{N}{2}(r - 2u\bar{\psi} + 3v\bar{\psi}^2)A^2 - \frac{T}{3}(u - 3v\bar{\psi})A^3 + \frac{vQ}{4}A^4,
\end{equation}
so that
\begin{equation}
u_{\mathrm{eff}} = u - 3v\bar{\psi}.
\end{equation}
The baseline tension directly tunes the triad coupling. It is not a nuisance parameter but the control parameter of the selection mechanism --- and it is precisely the quantity the ontology names as primitive.

\begin{figure}[h]
\centering
\fbox{\parbox{0.85\textwidth}{\centering [figure placeholder --- original Figure 4 not recovered]\\Phase diagram in the plane of baseline tension and quench depth. The full 24-mode $D_4$ standing wave occupies a contiguous band situated along the boundary with the disordered phase. Lower-symmetry sub-patterns occupy the higher-tension region.}}
\caption{Figure 4. Phase diagram: baseline tension versus quench depth.}
\end{figure}

\textbf{PREDICTION.} The full 24-mode pattern is the global minimum in a contiguous band occupying 11.85\% of the scanned plane, approximately $-0.33 \lesssim \bar{\psi} \lesssim 0.03$ depending on quench depth. The band is finite rather than measure-zero, and lies along the order/disorder boundary --- independently what the ``orders first'' criterion of \S II.3 predicts. Two calculations performed for different reasons agree on the location of the $D_4$ phase.

\textbf{Essential qualification, added in this version.} The phase diagram was computed with the bare quench depth, which tacitly assumes $q_0$ sits exactly on the $D_4$ shell. Any real mode on the shell incurs the penalty $(q_0^2 - |k|^2)^2$, so the effective quench depth is $r_{\mathrm{shell}} = r + (q_0^2 - 8)^2$. Solving the ordered branch symbolically at two values:
\begin{center}
\begin{tabular}{llll}
\toprule
$q_0^2$ & shell penalty & $r_{\mathrm{shell}}$ & ordered branch \\
\midrule
7.5 & 0.250 & $-0.100$ & none --- disordered state is the only minimum \\
8.0 & 0.000 & $-0.350$ & exists: $A_* = 0.0776$, $F = -0.00763$ \\
\bottomrule
\end{tabular}
\end{center}
\expansion{With $r = -0.35$: at $q_0^2 = 7.5$, penalty $(7.5-8)^2 = 0.25$, so $r_{\mathrm{shell}} = -0.35 + 0.25 = -0.10$; at $q_0^2 = 8.0$, penalty $0$, so $r_{\mathrm{shell}} = -0.35$. At a detuning of only 0.5 in $q_0^2$ the ordered branch vanishes entirely. The $D_4$ band is conditional on $q_0$ being tuned onto the shell. Since $a_0 = 2\pi/q_0$, this states that the substrate's resonant wavelength must match the lattice it forms --- self-consistent, but a condition rather than a free outcome. This qualification must accompany every citation of the 11.85\% figure.}
\flagbox{\textbf{Inconsistency preserved, not resolved:} the same section later states the $w_6 = 0$ band value as 11.90\%, against 11.85\% above --- a 0.05-point discrepancy in the source. Both figures are carried; the reader should treat the band fraction as $11.85$--$11.90\%$ pending the author's clarification. The phase-diagram percentages themselves are reported scan outcomes, \textbf{not re-run in this leg}.}

\textbf{Sixth-order terms.} Computed by FFT convolution and validated against brute-force third- and fourth-order values: for the $D_4$ shell $H = 911{,}040$. Adding $w_6$ narrows the band modestly (11.90\% at $w_6 = 0$ to 8.47\% at $w_6 = 0.2$) without destroying it. The $w_6 = 0$ value independently reproduces the 11.85\% obtained by different code.
\flagbox{The $H = 911{,}040$ count and the band percentages are reported computations, \textbf{not re-run in this leg}; the 11.85/11.90 tension above applies here too.}

\subsection{II.8 Stability of the Full Pattern}
\textbf{DERIVATION.} With $N = 24$, $T = 192$, $Q = 3384$, $r = -0.35$, $\bar{\psi} = -0.18$ and $q_0$ tuned, the stationarity condition yields a real ordered branch at per-mode amplitude $A_* = 0.0776$ (0.0775753 unrounded) with $F/V = -0.0076323 < 0$.

\textbf{Free-energy convention.} $F$ here is the excess over the uniform state, so $F = 0$ for the disordered phase. The absolute density differs by the uniform baseline $F/V|_{\mathrm{unif}} = -0.0034636$. Independent review reported a value of $-0.007728$ and attributed the difference to this baseline; that attribution is incorrect, since the baseline is $\sim 3.5 \times 10^{-3}$ while the discrepancy is $\sim 1 \times 10^{-4}$.

\textbf{The correct explanation, and why it strengthens the result.} The one-mode ansatz is variational: it is an upper bound on the true minimum, because a relaxed field develops harmonic content the ansatz excludes. Relaxing the one-mode state under conserved dynamics gives
\begin{center}
\begin{tabular}{ll}
\toprule
state & excess $F/V$ \\
\midrule
one-mode ansatz (analytic) & $-0.0076323$ \\
relaxed field (harmonics included) & $-0.0077276$ \\
independent review's value & $-0.007728$ \\
\bottomrule
\end{tabular}
\end{center}
The relaxed value reproduces the independently obtained number to five decimal places. The gap is harmonic lowering, not a convention mismatch, and it quantifies the accuracy of the one-mode approximation underpinning Part II: 1.25\%. That the Landau ansatz is accurate to better than 1.3\% against a full-field relaxation is a validation of the method, not a discrepancy in it.

\expansion{Gap arithmetic: $|-0.0077276 - (-0.0076323)| = 0.0000953 \approx 1\times10^{-4}$, against a baseline of $0.0034636 \approx 3.5\times10^{-3}$ --- the baseline is about 36 times larger than the gap, so it cannot explain the gap. Relative gap: $0.0000953/0.0077276 = 0.01233\ldots \approx 1.23\%$, quoted as 1.25\% (rounding of the underlying unrounded values). Rounding: $0.0775753 \to 0.0776$ at four decimal places. Verified in SymPy and python-flint.}

Seeding the substrate at exactly $A_* = 0.0776$ with all 24 modes in phase and evolving:
\begin{center}
\begin{tabular}{lllll}
\toprule
dynamics & final $\bar{\psi}$ & $D_4$ shell power & modes & rms modulation \\
\midrule
seed & $-0.18000$ & 1.0000 & 24/24 & 0.3802 \\
Model B (conserved) & $-0.18000$ & 1.0000 & 24/24 & 0.3806 \\
Model A (non-conserved) & $+0.00087$ & 1.0000 & 24/24 & 0.4206 \\
\bottomrule
\end{tabular}
\end{center}
Under conserved dynamics the baseline tension is held exactly, full occupancy is retained, and the amplitude is unchanged to four significant figures. The symmetric 24-mode $D_4$ standing wave is a genuine dynamical attractor, not merely a free-energy minimum. Under non-conserved dynamics the pattern also survives but the baseline tension drifts out of the stable band; its survival there is metastability. Every earlier simulation, initialized from noise, reached only a 6- or 12-mode sub-pattern. This is the first confirmation of the full symmetric state.
\flagbox{The evolution table and attractor claim are reported simulation outcomes, \textbf{not re-run in this leg}.}

\subsection{II.9 The Transition Is First Order}
At $\bar{\psi} = -0.18$ the effective quadratic coefficient is
\begin{equation}
r_{\mathrm{eff}} = r - 2u\bar{\psi} + 3v\bar{\psi}^2 = -0.35 + 0.36 + 0.097 = +0.107 > 0,
\end{equation}
so the disordered state is linearly stable and the $D_4$ state must be nucleated across a barrier, not grown from fluctuations.
\begin{center}
\begin{tabular}{lll}
\toprule
$\bar{\psi}$ & $r_{\mathrm{eff}}$ & character \\
\midrule
$-0.18$ & $+0.107$ & metastable --- requires nucleation \\
$-0.10$ & $-0.120$ & linearly unstable --- grows from noise \\
$0.00$ & $-0.350$ & linearly unstable \\
\bottomrule
\end{tabular}
\end{center}
\expansion{With $r = -0.35$, $u = v = 1$: at $\bar{\psi} = -0.18$: $-0.35 - 2(1)(-0.18) + 3(1)(0.0324) = -0.35 + 0.36 + 0.0972 = 0.1072 \approx +0.107$. At $\bar{\psi} = -0.10$: $-0.35 + 0.20 + 0.03 = -0.12$. At $\bar{\psi} = 0$: $-0.35$. All verified in SymPy and python-flint.}
\textbf{Consequence.} Any cosmological account within this framework must supply a nucleation mechanism: a seed, a fluctuation above critical size, or a path through the linearly unstable region at lower $|\bar{\psi}|$. Spontaneous emergence from a featureless substrate at this baseline tension does not occur. This was discovered by initially misreading a null simulation result and then correctly diagnosing it.

\subsection{II.10 Cross-Section to Three Dimensions}
\textbf{DERIVATION.} $D_4 \cap \{x_4 = 0\} = D_3 = \mathrm{FCC}$, with kissing number 12 --- the densest three-dimensional lattice packing~\cite{conway1999}. Every cross-section point satisfies the even-coordinate-sum condition defining FCC. The four-dimensional optimal lattice restricts to the three-dimensional optimal lattice.

\subsection{II.11 Scope of the Optimality Claim}
$D_4$ is not universally optimal. Independent verification via random-restart optimization across the Gaussian potential family --- which by Bernstein's theorem~\cite{bernstein1929} spans the completely-monotonic class --- recovers $D_4$ as the exact optimum (to $\sim 10^{-13}$ relative) across $\beta \in [0.05, 2.4] \cup [3.1, 12]$, but a lower-symmetry competitor wins in a narrow window $\beta \in [2.6, 3.0]$ by at most 0.007\%. This independently confirms the published non-universal-optimality result~\cite{cohn2007} by direct discovery; see also the recent optimality and uniqueness analysis~\cite{delaat2024}. The universal-optimality framework itself --- configurations minimizing energy over all completely monotonic potentials --- is due to Cohn and Kumar~\cite{cohn2007b}. The defensible statement is that $D_4$ minimizes the specific triad-reinforcing functional relevant to standing-wave selection, and is the unique optimal kissing configuration in four dimensions~\cite{musin2008}, but does not minimize every completely-monotonic potential.
\flagbox{The optimization scan ($\beta$ windows, $10^{-13}$, 0.007\%) is a reported computation, \textbf{not re-run in this leg}.}

% ============================ PART III ============================
\section{PART III --- DYNAMICS}

\subsection{III.1 The Discrete Action (Legacy Formulation)}
For the site-field formulation, retained here with scope stated:
\begin{equation}
S = \sum_{n,s} a_0^4 \Delta t \left[ \frac{M_*}{2}\left(\frac{\phi(n,s+1) - \phi(n,s)}{\Delta t}\right)^2 - \frac{J}{2a_0^2}\sum_{\alpha \in \Phi_{D_4}} (\phi(n+\alpha,s) - \phi(n,s))^2 - \frac{\mu^2}{2}\phi(n,s)^2 \right].
\end{equation}

\textbf{DERIVATION --- equation of motion, by symbolic variation with all 24 neighbour contributions entered explicitly:}
\begin{equation}
M_*\left[\phi(n,s+1) - 2\phi(n,s) + \phi(n,s-1)\right] = \frac{2J\Delta t^2}{a_0^2}\sum_{\alpha} \left[\phi(n+\alpha,s) - \phi(n,s)\right] - \Delta t^2\mu^2\phi(n,s).
\end{equation}

\textbf{DERIVATION --- exact dispersion relation.} Plane-wave substitution reduced symbolically gives the closed form
\begin{equation}
\sum_{\alpha \in \Phi_{D_4}} \left[\cos(a_0 k\cdot\alpha) - 1\right] = 4\sum_{1\leq i<j\leq 4}\cos(a_0k_i)\cos(a_0k_j) - 24,
\end{equation}
\begin{equation}
\sin^2\left(\frac{\omega\Delta t}{2}\right) = \frac{\Delta t^2\,2J}{4M_*a_0^2}\left[\sum_{\alpha}\left(1 - \cos(a_0k\cdot\alpha)\right) + \mu^2\right].
\end{equation}
\expansion{The dispersion identity holds pair by pair: for fixed $(i,j)$ the four roots give $\cos(a(k_i+k_j)) + \cos(a(k_i-k_j)) + \cos(a(-k_i+k_j)) + \cos(a(-k_i-k_j)) = 4\cos(ak_i)\cos(ak_j)$ by the cosine addition formula ($\cos(X+Y) + \cos(X-Y) = 2\cos X\cos Y$, applied twice). Summing over the 6 pairs and subtracting 24 (one per root) gives the stated identity. SymPy 1.14.0 confirms the pair identity symbolically and the full identity at 20 random numeric points to $< 10^{-9}$.}
\textbf{Cross-check.} Taylor-expanding reproduces the isotropic result of \S II.1 term for term, via a completely independent route --- a direct root-sum versus a full Euler--Lagrange variation. Agreement between independent derivations is evidence, not merely internal consistency.

\textbf{DERIVATION --- Hamiltonian positivity.} The energy is a sum of three non-negative terms given $M_*, J > 0$ and $\mu^2 \geq 0$: the vacuum is stable by construction, with no ghosts or runaway modes.

\subsection{III.2 Stability Bound}
\textbf{DERIVATION.} Real solutions require $\sin^2(\omega\Delta t/2) \leq 1$ for every $k$. The maximum of the spatial sum over the Brillouin zone was found analytically and confirmed by global numerical optimization: exactly 32, at $a_0k = (0,0,\pi,\pi)$-type points. Matching the small-$k$ limit to a continuum wave speed fixes $J = M_*c^2/12$, giving
\begin{equation}
\Delta t \leq \Delta t_{\max} = \frac{\sqrt{3}}{2}\frac{a_0}{c} \approx 0.866\,\frac{a_0}{c}.
\end{equation}

\begin{figure}[h]
\centering
\fbox{\parbox{0.85\textwidth}{\centering [figure placeholder --- original Figure 5 not recovered]\\The stability quantity at the worst-case Brillouin-zone mode against time step.}}
\caption{Figure 5. Stability quantity at the worst-case mode.}
\end{figure}

\textbf{EXCLUDED --- both time-step values stated in v92.}
\begin{center}
\begin{tabular}{lll}
\toprule
candidate $\Delta t$ & $\sin^2(\omega\Delta t/2)$ at worst-case $k$ & verdict \\
\midrule
$a_0/c$ & $4/3 \approx 1.333$ & unstable \\
$2\pi a_0/c$ & $16\pi^2/3 \approx 52.6$ & badly unstable \\
$(\sqrt{3}/2)a_0/c$ (derived) & 1 (exactly marginal) & stable \\
\bottomrule
\end{tabular}
\end{center}
\expansion{With $J = M_*c^2/12$ and $\mu = 0$, the prefactor in (7) is $\Delta t^2\cdot 2J/(4M_*a_0^2) = \Delta t^2c^2/(24a_0^2)$; at the worst case the bracket equals 32, so $\sin^2 = \tfrac{4}{3}(\Delta t\,c/a_0)^2$. Hence: $\Delta t = a_0/c \Rightarrow 4/3 = 1.333\ldots$; $\Delta t = 2\pi a_0/c \Rightarrow \tfrac{4}{3}\cdot 4\pi^2 = 16\pi^2/3 = 52.6378\ldots \approx 52.6$; $\Delta t = \tfrac{\sqrt{3}}{2}a_0/c \Rightarrow \tfrac{4}{3}\cdot\tfrac{3}{4} = 1$ exactly. The $a_0/c$ candidate exceeds the bound by $1/(\sqrt{3}/2) = 2/\sqrt{3} = 1.1547\ldots$, i.e.\ roughly 15\%. All verified in SymPy and python-flint. The ``maximum exactly 32'' itself is a stated analytic/numerical result, \textbf{not re-derived in this leg}.}
Neither of v92's stated time steps is dynamically stable once the coupling is honestly calibrated. This is sharper than the previously identified factor-of-$2\pi$ ambiguity: even the smaller candidate exceeds the limit by roughly 15\%.

\subsection{III.3 The Non-Relativistic Limit}
\textbf{Diagnosis of the v92 failure.} The manuscript constructed a carrier $e^{-i\Omega_P sT_P}$ with $T_P = 2\pi/\Omega_P$, so the carrier phase per step is exactly $2\pi$ and $e^{-i\theta s} = 1$ for every integer $s$. This is not merely insufficient oscillation: the positive- and negative-frequency basis functions become identical, so the envelope decomposition is mathematically degenerate. The subsequent small-angle expansion of $\sin(2\pi)$ compounds the error.

\textbf{DERIVATION --- corrected envelope equation.} With a genuine carrier ($\theta \neq 0, \pi$) and ansatz $\phi = \Psi e^{-i\theta s} + \mathrm{c.c.}$, symbolic Taylor expansion with the rest-mass term exactly absorbed via $\sin^2(\theta/2) = \Delta t^2\mu^2/4M_*$:
\begin{equation}
i\dot{\Psi} = -\frac{\Delta t\,J}{M_*a_0^2\sin\theta}\sum_{\alpha \in \Phi_{D_4}} \left[\Psi(n+\alpha,s) - \Psi(n,s)\right],
\end{equation}
with continuum effective mass
\begin{equation}
m_{\mathrm{eff}} = \frac{M_*a_0\sin\theta}{c\,\Delta t}.
\end{equation}
\textbf{PREDICTION --- quantified validity regime.} Comparing the exact transcendental dispersion against the quadratic approximation mode by mode:
\begin{center}
\begin{tabular}{ll}
\toprule
$a_0|k|$ & relative frequency error \\
\midrule
0.02 & $4 \times 10^{-5}$ \\
0.10 & 0.1\% \\
0.30 & 1.0\% \\
1.00 & 11\% \\
3.14 (zone edge) & 133\% \\
\bottomrule
\end{tabular}
\end{center}
\begin{figure}[h]
\centering
\fbox{\parbox{0.85\textwidth}{\centering [figure placeholder --- original Figure 6 not recovered]\\Exact dispersion against the envelope approximation, and the resulting relative error.}}
\caption{Figure 6. Exact dispersion versus envelope approximation.}
\end{figure}
The approximation is accurate to well under 1\% for $a_0|k| \lesssim 0.3$, degrading smoothly beyond. This is a quantified boundary rather than an assumed one.
\flagbox{The validity table is a reported mode-by-mode comparison, \textbf{not re-run in this leg} (no dispersion-comparison code in the source text).}

\subsection{III.4 Why Four Dimensions}
\textbf{DERIVATION.} Under $\psi(x) \to \psi(\lambda x)$ in $d$ dimensions the terms of a scalar energy scale as
\begin{equation}
\int |\nabla\psi|^2 \sim \lambda^{2-d}, \qquad \int V(\psi) \sim \lambda^{-d}, \qquad \int (\nabla^2\psi)^2 \sim \lambda^{4-d}.
\end{equation}
Stationarity of the two-derivative theory at $\lambda = 1$ requires $(d-2)I_{\mathrm{grad}} + d\,I_{\mathrm{pot}} = 0$, which for $d \geq 2$ forces $I_{\mathrm{pot}} \leq 0$. With a positive-definite potential there is no stationary solution, and where one exists the second variation is negative. Derrick's theorem~\cite{derrick1964} forbids stable Kerr solitons in four dimensions. But the four-derivative Swift--Hohenberg term is exactly marginal at $d = 4$.
\expansion{Scaling step by step: $d^dx \to \lambda^{-d}d^dx$ under $x \to \lambda x$; $\nabla \to \lambda\nabla$ so $|\nabla\psi|^2 \to \lambda^2|\nabla\psi|^2$, giving $\lambda^{2-d}$; $V(\psi)$ carries no derivatives, giving $\lambda^{-d}$; $(\nabla^2\psi)^2 \to \lambda^4(\nabla^2\psi)^2$, giving $\lambda^{4-d}$. Marginality means exponent zero: $4 - d = 0$, i.e.\ $d = 4$ uniquely. Derrick: $\tfrac{d}{d\lambda}[\lambda^{2-d}I_{\mathrm{grad}} + \lambda^{-d}I_{\mathrm{pot}}]_{\lambda=1} = (2-d)I_{\mathrm{grad}} - d\,I_{\mathrm{pot}} = 0$, i.e.\ $(d-2)I_{\mathrm{grad}} + d\,I_{\mathrm{pot}} = 0$. Verified symbolically in SymPy.}
\begin{figure}[h]
\centering
\fbox{\parbox{0.85\textwidth}{\centering [figure placeholder --- original Figure 7 not recovered]\\Scaling exponents against spatial dimension. The four-derivative term is scale-invariant precisely at $d = 4$, where the two-derivative and potential terms are both strongly contracting.}}
\caption{Figure 7. Scaling exponents versus dimension.}
\end{figure}
This is a new structural argument for four-dimensionality, independent of packing and kissing-number considerations. Four is the unique dimension in which the curvature term of a Swift--Hohenberg substrate neither dominates nor vanishes under rescaling, so localized structure is driven neither to collapse nor to dispersal at leading order. The dimension hosting the optimal lattice~\cite{korkine1877} is also the dimension in which the pattern-forming operator is marginal.

% ============================ PART IV ============================
\section{PART IV --- ALGEBRAIC STRUCTURE}

\subsection{IV.1 The $SO(8)$ Exclusion Theorem}
\textbf{EXCLUDED --- $\mathfrak{so}(8)$ cannot be the gauge algebra. Two independent proofs.}

\textbf{Chirality.} A representation is self-conjugate precisely when $-1$ lies in the Weyl group~\cite{humphreys1972}; for $D_n$ this holds if and only if $n$ is even. Computed by two structurally different algorithms:
\begin{center}
\begin{tabular}{lllll}
\toprule
method & $|W(D_4)|$ & $-I \in W$? & $|W(D_5)|$ & $-I \in W$? \\
\midrule
signed-permutation construction & 192 & yes & 1920 & no \\
generation from simple reflections & 192 & yes & 1920 & no \\
\bottomrule
\end{tabular}
\end{center}
\expansion{Weyl orders: $|W(D_n)| = 2^{n-1}n!$, so $|W(D_4)| = 2^3\cdot 24 = 192$ and $|W(D_5)| = 2^4\cdot 120 = 1920$. Verified in SymPy. The two-algorithm agreement is a stated result, \textbf{not re-run in this leg}.}
Since $-1 \in W(D_4)$, every representation of $\mathfrak{so}(8)$ is self-conjugate. The three eight-dimensional representations are all real and permuted by triality. An algebra with only real representations is vector-like and cannot produce chiral fermion couplings. The Standard Model is chiral. The contrast is exact and instructive: $-1 \notin W(D_5)$, which is precisely why $\mathfrak{so}(10)$ admits the complex 16 and functions as a viable chiral unified group.

\textbf{Embedding.} Independently, the Standard Model algebra has rank 4, equal to that of $\mathfrak{so}(8)$, so it could only appear as a maximal-rank subalgebra. By Borel--de Siebenthal~\cite{borel1949} these are $\mathfrak{su}(4) \oplus \mathfrak{u}(1)$, $\mathfrak{su}(2)^{\oplus 4}$, and $\mathfrak{su}(2)^{\oplus 3} \oplus \mathfrak{u}(1)$. None contains $\mathfrak{su}(3) \oplus \mathfrak{su}(2)$: the $\mathfrak{su}(2)$ towers cannot contain simple rank-2 $\mathfrak{su}(3)$, and $\mathfrak{su}(3) \oplus \mathfrak{su}(2) \subset \mathfrak{su}(4)$ fails because faithful actions on commuting subspaces require at least five complex dimensions while $\mathfrak{su}(4)$ acts on $\mathbb{C}^4$.
\begin{figure}[h]
\centering
\fbox{\parbox{0.85\textwidth}{\centering [figure placeholder --- original Figure 8 not recovered]\\Confirmed regular-subalgebra branchings of $\mathfrak{so}(8)$, showing that neither route reaches the Standard Model jointly.}}
\caption{Figure 8. Regular-subalgebra branchings of $\mathfrak{so}(8)$.}
\end{figure}
\textbf{Also excluded:} Pati--Salam~\cite{pati1974} (rank 5) cannot be reached from $\mathfrak{so}(8)$ (rank 4) by any route whatsoever. The v92 cascade $\mathfrak{so}(8) \to \mathfrak{g}_2 \oplus \mathfrak{u}(1) \to$ Pati--Salam additionally increases dimension from 15 to 21, which no subalgebra chain can do.
\expansion{Dimension arithmetic, clarified in this rebuild: $\dim\mathfrak{g}_2 = 14$, so $\dim(\mathfrak{g}_2 \oplus \mathfrak{u}(1)) = 14 + 1 = 15$; $\dim(\text{Pati--Salam}) = \dim\mathfrak{su}(4) + \dim\mathfrak{su}(2) + \dim\mathfrak{su}(2) = 15 + 3 + 3 = 21$; ranks $3 + 1 + 1 = 5 > 4 = \mathrm{rank}\,\mathfrak{so}(8)$. The source's ``$15 \to 21$'' is the cascade's intermediate-to-final step, not $\dim\mathfrak{so}(8) = 28$.}

\subsection{IV.2 The Forcing Result}
\textbf{DERIVATION --- the phase sector has no internal degrees of freedom.} A real $D_4$ density wave has 24 modes forming 12 independent $\pm$ pairs, hence 12 phases, constrained by closed triads $\phi_a + \phi_b + \phi_c = 0$.
\begin{center}
\begin{tabular}{ll}
\toprule
quantity & value \\
\midrule
independent phases & 12 \\
unordered closed triads & 32 \\
rank of constraint matrix & 8 \\
remaining degrees of freedom & 4 \\
translation subspace dimension & 4 \\
leftover non-translational DOF & 0 \\
\bottomrule
\end{tabular}
\end{center}
\expansion{$12 - 8 = 4$ remaining degrees of freedom; $4 - 4 = 0$ leftover after translations. The inputs (12 phases, 32 unordered triads, rank 8) are stated counts; the subtractions are the verified arithmetic. Note the distinction kept throughout: 192 counts \emph{ordered} triads (\S II.5), 32 counts \emph{unordered} closed triads here --- different countings of related objects, not a contradiction.}
Every phase is fixed by triad locking up to a rigid four-dimensional translation. The order-parameter space is the torus $\mathbb{R}^4/\Lambda$, carrying no internal continuous symmetry.

\textbf{DERIVATION --- the homotopy is equally restrictive.} For $M = T^4$~\cite{mermin1979}: $\pi_1 = \mathbb{Z}^4$, $\pi_2 = 0$, $\pi_3 = 0$. Dislocations exist; monopoles do not; skyrmions and hopfions do not.

\textbf{Consequence.} A single real scalar substrate can produce the $D_4$ pattern and dislocations, but cannot produce gauge fields or point-like defects. Internal structure is forced, not chosen. The question was never which lattice symmetry to promote to a gauge group, but what internal space the substrate carries.

\subsection{IV.3 The Substrate Identifies Its Own Internal Space}
\textbf{DERIVATION --- the 24-cell is a group.} Exhaustive evaluation of all 576 products of the 24 vertices in unit normalization --- the eight axis quaternions $\pm1, \pm i, \pm j, \pm k$ (the Hurwitz integer units~\cite{hurwitz1896}) and the sixteen $(\pm1 \pm i \pm j \pm k)/2$ --- confirms closure under quaternion multiplication, with identity and inverses. The 24-cell vertices form the binary tetrahedral group $2T$~\cite{coxeter1957}, a finite subgroup of $SU(2) \cong S^3$.
\expansion{Checked in this rebuild by exhaustive Python enumeration (second tool): 24 vertices; all $24^2 = 576$ products land in the set; the 8 axis quaternions are closed; $Q_8$ is normal in $2T$ (all conjugates stay in $Q_8$).}
This is the pivot of the gauge resolution. The $D_4$ pattern is not merely a shape with a symmetry group; its vertex set is a group, and that group sits inside $SU(2)$. The internal space is therefore identified by the pattern rather than posited. It also repairs the homotopy deficit of \S IV.2: $\pi_3(S^3) = \mathbb{Z}$.

\textbf{DERIVATION --- coset structure.} The eight axis quaternions form $Q_8$, verified closed and normal in $2T$, with
\begin{equation}
2T/Q_8 \cong \mathbb{Z}_3, \qquad \text{three cosets of eight}.
\end{equation}
The $8+8+8$ splitting of the 24-cell is the coset decomposition of an actual group, not a labelling convention.
\expansion{Enumeration confirms: 3 cosets, each of size 8, partitioning the 24 elements.}

\textbf{DERIVATION --- McKay correspondent.} Computing genuine conjugacy classes by conjugation orbits (trace and element order are too coarse and merge classes) gives seven classes of sizes 4,4,4,4,1,1,6 summing to 24. Requiring dimension-squares to sum to the group order with exactly three one-dimensional representations yields a unique solution:
\begin{equation}
\dim(\mathrm{irreps\ of\ } 2T) = \{1, 1, 1, 2, 2, 2, 3\}.
\end{equation}
\begin{center}
\begin{tabular}{lllll}
\toprule
affine algebra & marks & nodes & match \\
\midrule
$E_6$ & 1,1,1,2,2,2,3 & 7 & yes \\
$E_7$ & 1,1,2,2,2,3,3,4 & 8 & no \\
$E_8$ & 1,2,2,3,3,4,4,5,6 & 9 & no \\
\bottomrule
\end{tabular}
\end{center}
\expansion{Class sizes: $4+4+4+4+1+1+6 = 24$; class count 7 equals the irrep count. Dimension squares: $1+1+1+4+4+4+9 = 24 = |2T|$. Both confirmed by exhaustive enumeration (conjugacy orbits) in this rebuild. The marks tables are the standard affine-Dynkin marks, carried as stated.}
$2T$ corresponds to affine $E_6$~\cite{mckay1980}. An earlier circulated audit document asserted correspondence to affine $E_8$ and built a sterile-neutrino prediction upon it; affine $E_8$ corresponds to the binary icosahedral group of order 120, and any prediction resting on that identification does not follow.

\textbf{DERIVATION --- $E_6$ is chiral.} Computing the longest Weyl element identically for all three algebras from their Cartan matrices:
\begin{center}
\begin{tabular}{llll}
\toprule
algebra & rank & $w_0 = -1$? & verdict \\
\midrule
$D_4 = \mathfrak{so}(8)$ & 4 & yes & non-chiral \\
$D_5 = \mathfrak{so}(10)$ & 5 & no & chiral \\
$E_6$ & 6 & no & chiral \\
\bottomrule
\end{tabular}
\end{center}
$E_6$ passes the identical test $\mathfrak{so}(8)$ fails. It has the complex 27~\cite{slansky1981}, and the rank chain $E_6(6) \supset SO(10)(5) \supset SU(5)(4) \supset \mathrm{SM}(4)$ is non-increasing --- the standard $E_6$~\cite{gursey1976} via $SO(10)$~\cite{fritzsch1975} and $SU(5)$~\cite{georgi1974} breaking route. One generation of Standard Model fermions embeds in a single 27.
\expansion{The $E_6$ Cartan determinant is 3 (SymPy and Singular agree), and its Smith normal form is $\mathrm{diag}(1,1,1,1,1,3)$ (SymPy), giving $P/Q \cong \mathbb{Z}_3$ as used in \S V.2. The $w_0$ table itself is a stated Lie-algebra computation, carried as stated.}

\begin{figure}[h]
\centering
\fbox{\parbox{0.85\textwidth}{\centering [figure placeholder --- original Figure 9 not recovered]\\Left: the 24 vertices coloured by coset of $Q_8$ in $2T$. Right: the resulting structural chain from substrate to Standard Model.}}
\caption{Figure 9. Cosets of $Q_8$ in $2T$; substrate-to-Standard-Model chain.}
\end{figure}

\subsection{IV.4 The Gauge Decision}
The gauge algebra of IRH is $E_6$, arising from the substrate's internal space $S^3 \cong SU(2)$ via the binary tetrahedral group and the McKay correspondence. What makes this defensible rather than convenient is \S IV.3. $E_6$ is not selected for being large enough to contain what is needed. The pattern established on independent dynamical grounds turns out to be a group whose McKay correspondent is $E_6$. The gauge structure is read off the substrate, not fitted to the target.

\textbf{What is given up.} Triality is no longer the generation mechanism and $\mathfrak{so}(8)$ is no longer a gauge algebra. Triality survives as a symmetry of the defect classification --- the $D_4$ lattice point group is $W(F_4)$ of order 1152~\cite{conway1999}, which is $W(D_4)$ of order 192 extended by the $S_3$ triality automorphism --- but it is no longer load-bearing.
\expansion{$|W(F_4)| = 1152 = 2^7\cdot 3^2$ and $1152/192 = 6 = |S_3|$, consistent with the stated extension. Verified arithmetically in SymPy.}

% ============================ PART V ============================
\section{PART V --- MATTER AND DEFECTS}

\subsection{V.1 The Defect Spectrum}
With internal space $S^3$, the order-parameter space is $M = T^4 \times S^3$, giving $\pi_1(M) = \mathbb{Z}^4$, $\pi_2(M) = 0$, $\pi_3(M) = \mathbb{Z}$. In four spatial dimensions a defect of codimension $k$ is classified by $\pi_{k-1}$~\cite{mermin1979}:
\begin{center}
\begin{tabular}{lll}
\toprule
defect & classified by & result \\
\midrule
codimension 2 (sheets) & $\pi_1 = \mathbb{Z}^4$ & dislocation sheets, Burgers vectors in $D_4$ \\
codimension 3 (lines) & $\pi_2 = 0$ & none \\
codimension 4 (points) & $\pi_3 = \mathbb{Z}$ & point-like solitons, integer charge \\
\bottomrule
\end{tabular}
\end{center}
The point-like, integer-charged topological objects required by the standing-wave ontology do exist once the internal space is $S^3$. They did not exist for a single real scalar. The gauge decision repaired this as a side effect rather than by design. Elementary dislocation charges: 24, the shortest $D_4$ lattice vectors. Point group: $W(F_4)$, order 1152.

\subsection{V.2 Generations: A Negative Result}
The question was posed sharply in v92.6 and declined there as unproven: does the $\mathbb{Z}_3$ from $2T$ produce three copies of the 27, or are the two threes unrelated? Five candidate mechanisms were tested.

\textbf{The coset $\mathbb{Z}_3$ is a diagram automorphism.} The three one-dimensional irreps of $2T$ are the three mark-1 tip nodes of affine $E_6$. The diagram automorphism group has order 6, permuting the arms. An automorphism permutes nodes; it does not generate multiplicity.

\textbf{The centre $\mathbb{Z}_3$ grades, it does not count.} Smith normal form~\cite{newman1972} of the $E_6$ Cartan matrix gives invariant factors $(1,1,1,1,1,3)$, so $P/Q \cong \mathbb{Z}_3$. The 27 carries class 1, the $\overline{27}$ class 2, the adjoint 78 class 0. The centre distinguishes a representation from its conjugate and from the adjoint. Grading three different representations is not producing three copies of one.
\expansion{Smith normal form $\mathrm{diag}(1,1,1,1,1,3)$ and $\det = 3$ verified in SymPy (see \S IV.3 expansion). The class assignments (27: 1, $\overline{27}$: 2, 78: 0) are standard representation theory, carried as stated.}

\textbf{Anomaly cancellation imposes nothing.} $E_6$ has exponents $\{1,4,5,7,8,11\}$ --- no exponent 2, hence no third-order Casimir. The cubic anomaly coefficient vanishes identically for every representation~\cite{slansky1981}, so any number of 27's is consistent.

\textbf{The index route is closed twice over.} Generation number in compactifications is an index, not a symmetry. On $S^3/2T$ the Dirac index vanishes by Lichnerowicz vanishing~\cite{lichnerowicz1963} (positive scalar curvature) and, more basically, on any closed odd-dimensional manifold, the index density being built from Pontryagin classes concentrated in degrees divisible by four~\cite{lawson1989}.

\textbf{VERDICT --- three generations is an INPUT, not an output.} Three distinct objects were being conflated: a symmetry permuting cosets, a symmetry grading representations, and a multiplicity in the field content. The first two are real and structurally elegant. Neither is the third. This must be stated plainly because the coincidence is seductive. The 24-cell really does decompose into three octets; $E_6$ generations really do come in threes; the numbers really do match. The match is not a mechanism. The framework does not predict three generations and has not done so since $\mathfrak{so}(8)$ was excluded.
\gradebox{Three generations: INPUT, grade unchanged by this rebuild. Any future text that presents generation count as derived from IRH would be an overstatement.}

\subsection{V.3 An Inherited Liability: Monopoles}
Any simple gauge group broken to a subgroup containing an unbroken $U(1)$ has $\pi_2(G/H) = \mathbb{Z}$ and therefore stable magnetic monopoles~\cite{preskill1984}. $E_6 \to \cdots \to SU(3)\times SU(2)\times U(1)$ is of exactly this form.

\textbf{EXCLUDED --- monopole absence, claimed in earlier versions, is inconsistent with an $E_6$ gauge group.} IRH inherits the standard unified-theory monopole problem and requires the standard resolution: inflationary dilution, or a breaking pattern avoiding the relevant topology.

% ============================ PART VI ============================
\section{PART VI --- FALSIFIABLE PREDICTIONS}
The framework makes the following predictions, each independent of the data it would be tested against.

\textbf{VI.1 Triad phase locking.} For a real field with $u > 0$, the cubic term is minimized when $\phi_1 + \phi_2 + \phi_3 = 0 \pmod{2\pi}$ on every closed triad. This is a hard signature distinguishing a genuine $D_4$ standing wave from a coincidental superposition of modes on the right shell. Measured in simulation at an order parameter of 0.26 against a maximum of 1 --- real but incomplete locking, consistent with a defected state in a small box. Test: larger volume should drive this toward unity.
\flagbox{The 0.26 order parameter is a reported simulation measurement, \textbf{not re-run in this leg}.}

\textbf{VI.2 The baseline-tension band.} The full 24-mode pattern is the global minimum only for $\bar{\psi}$ in a definite band along the order/disorder boundary, and only for $q_0$ tuned onto the $D_4$ shell. Outside either condition, lower-symmetry patterns or no pattern result.

\textbf{VI.3 Validity boundary of the non-relativistic limit.} Under 1\% frequency error for $a_0|k| \lesssim 0.3$, degrading as tabulated in \S III.3.

\textbf{VI.4 Stability bound.} $\Delta t \leq (\sqrt{3}/2)a_0/c$ for any discrete implementation.

\textbf{VI.5 Isotropy to fourth order.} No anisotropic quartic Lorentz-violating term arises from nearest-neighbour $D_4$ coupling. Anisotropy enters at sixth order or via next-nearest-neighbour couplings. This is a prediction of absence and is falsified by any observation of quartic-order directional dependence attributable to this mechanism.

\textbf{VI.6 Defect spectrum.} Dislocation sheets with Burgers vectors in $D_4$, integer-charged point solitons, and no line defects.

\textbf{VI.7 Monopoles.} Predicted, not excluded --- see \S V.3.

\section{PART VI-A --- INDEPENDENT VERIFICATION}
An independent party re-derived the load-bearing claims of this framework from scratch, including a Lean~4 formalization. The results are recorded here with the same standard applied to everything else: confirmations were re-checked rather than accepted, and two of the reviewer's own findings did not survive that check.

\subsection{VI-A.1 Confirmed}
The lattice mathematics, the exact isotropy moments, the resonance maximum of 24, the CFL derivation, the dispersion identity, the $D_4 \to$ FCC cross-section, the Weyl group orders and the $-1 \in W$ distinction, the quaternion group checks, and the $u_{\mathrm{eff}}$ identity were all independently reproduced.

The attractor was independently confirmed. Seeding all 24 modes in phase at $A_* = 0.0776$, the state was found stationary to $\sim 10^{-5}$ under exact convolutions, with Model B holding the amplitude to four significant figures at $D_4$ power 1.0000, and Model A's drift to $\bar{\psi} = +0.00087$ with rms 0.4206 matching the values reported in \S II.8 exactly.

Selection dynamics reproduced. The \S II.6 fair test was reproduced at mechanism level by an independent implementation: with triads active, $D_4$ captured 77.1\% of spectral power on direct quench (author's run: 84\%), rising to 89.5\% under 14000 annealing steps and still increasing (author's run: 100\%); with triads removed, $D_4$ fell to 21.1\% against 78.9\% for the competitor (author's run: 44\% against 56\%). Exact shares depend on seed, scheme, and annealing protocol; the causal claim --- the cubic term flips the winner --- reproduces unambiguously.
\flagbox{Reproduction percentages are reported outcomes of the independent implementation, \textbf{not re-run in this leg}.}

A sharpening. The reviewer noted that the selection quantities are exact rationals and simple radicals rather than decimals. This is correct and has been incorporated into \S II.5. It is a genuine strengthening: a mechanism producing $64/141$, $16/45$, $73/423$ and an identity $\tau_{D_4}/\tau_{\mathrm{BCC}} = \sqrt{2}$ \textbf{[corrected]} is exhibiting structure, not fitted parameters.

\subsection{VI-A.2 Formal verification status}
The accompanying Lean~4 development~\cite{demoura2015} contains no sorry, no axiom, no admit, and no native\_decide, verified by direct search of every file. The theorem statements are honest about what they prove: \texttt{d4\_fourth\_moment} states the identity for a general wavevector $k : \mathrm{Fin}\,4 \to \mathbb{R}$, and the 24 listed roots are the correct root system. Nine theorems are covered, spanning isotropy, the dispersion identity, the CFL bound, the selection criterion and its exact rationals, the $u_{\mathrm{eff}}$ identity, the cross-section, and $-1 \in W(D_4)$.

Build status (v93.1). The Lean development was compiled end to end under Lean~4 with Mathlib~\cite{mathlib2020}: \texttt{lake build} exited 0 (``Build completed successfully (3101 jobs)''), producing all eight module .olean files. The claim is therefore ``machine-checked end to end,'' not merely ``formally stated and independently re-derived.'' The v93.0 stated limitation --- that the development had been audited by inspection but not compiled --- is superseded.
\flagbox{The Lean build claim is carried as stated in v93.1. No Lean build was run in this leg (that is another leg's job, per the brief); this rebuild neither confirms nor re-grades it.}

\subsection{VI-A.3 Two reviewer findings that did not survive checking}
\textbf{The free-energy gap.} The reviewer attributed a 1.2\% difference to the uniform baseline convention. That attribution is arithmetically wrong: the baseline is $3.5\times10^{-3}$ while the gap is $1\times10^{-4}$. The correct explanation is harmonic content beyond the one-mode ansatz, and it reproduces the reviewer's own number to five decimals (\S II.8). The result is stronger than the reviewer's reading of it: it quantifies the one-mode approximation's accuracy at 1.25\%.
\expansion{Gap $0.0000953$ versus baseline $0.0034636$: see the \S II.8 expansion. The ``1.2\%'' figure is the reviewer's rounded description of the same $\approx$1.23\% relative gap.}

\textbf{The periodic-image count.} The reviewer stated that the resonance maximum recurs 16 times on a periodic grid. The correct count is 32. The factor of two is the $D_4$ checkerboard structure: corner and body-centre sublattices both satisfy the coincidence condition. The reviewer's underlying point --- that ``unique'' must be qualified as ``modulo lattice periodicity'' --- stands and has been incorporated.
\expansion{$2^4 = 16$ repeats $\times$ 2 sublattices $= 32$: see the \S II.2 expansion.}

\subsection{VI-A.4 Numerical scheme}
The reviewer reported a late-onset instability in their semi-implicit Model B integration at $t \gtrsim 12$. The scheme used here was re-run over the same horizon and showed none: rms modulation held at 0.380565 and $\max|\psi|$ at 1.689598, constant to six significant figures from $t = 2$ through $t = 18$. That the seeded state is stationary to this precision over a long horizon is an additional confirmation that it is a genuine fixed point rather than a slowly decaying transient. The difference is implementation-specific; dealiasing and timestep choice for the conserved scheme are documented in Appendix B.
\flagbox{The long-horizon stationarity figures are reported re-run outcomes, \textbf{not re-run in this leg}.}

% ============================ PART VII ============================
\section{PART VII --- RETRACTIONS FROM v92}
The following claims in v92 are withdrawn. Each is withdrawn for a stated reason, not for caution.
\begin{center}
\begin{tabular}{lll}
\toprule
v92 claim & status & reason \\
\midrule
$\xi_2 = -1/384$ anisotropy coefficient & retracted & $D_4$ fourth moment is exactly isotropic; the value derives from a hypercubic stand-in, and the stated substitution yields $-1/768$, off by a factor of 2 \\
$\mathfrak{so}(8) \to \mathfrak{g}_2 \to$ Pati--Salam cascade & retracted & dimension increases $15 \to 21$ and rank $3 \to 5$; impossible [see \S IV.1 clarification: $15 = \dim(\mathfrak{g}_2\oplus\mathfrak{u}(1))$, $21 = \dim$ Pati--Salam; rank $4 \to 5$ for the full chain] \\
$SO(8)$ as gauge group & retracted & non-chiral: $-1 \in W(D_4)$ \\
Three generations from triality & retracted & five mechanisms closed; now an acknowledged input \\
$\Delta t = a_0/c$ and $\Delta t = T_P$ & retracted & both violate the derived CFL bound \\
Schr\"odinger emergence & corrected & original carrier degenerate; replaced by \S III.3 \\
Lorentzian signature from $i^2 = -1$ & retracted & canonical $i\hbar$ is present in Euclidean field theory also; does not select signature \\
Einstein equations from elasticity & retracted & analogy, not derivation; see \S IX.3 \\
Monopole absence & retracted & inconsistent with $E_6$ breaking topology \\
Higgs VEV $v_0$, $A_3$ mass and widths, CKM & retracted & computed in an excluded framework, with independent arithmetic errors \\
``Zero empirical calibration'', ``100\% closure'' & retracted & not defensible in any version \\
$\alpha^{-1} = 137 + 1/(28 - \pi/14)$ as derived & downgraded & arithmetic verified exact; physical derivation absent. Numerical ansatz only \\
$R = 2589$ normalization & downgraded & arithmetic exact; the subtraction of 3 is not forced by any displayed symmetry \\
\bottomrule
\end{tabular}
\end{center}
\expansion{Retraction arithmetic checked in this rebuild: $(-1/384)/(-1/768) = 2$, confirming the factor-of-2 error. The $\alpha^{-1}$ ansatz: $\pi/14 = 0.224399\ldots$, $28 - \pi/14 = 27.775600\ldots$, $1/27.775600\ldots = 0.0360028\ldots$, plus 137 gives $137.0360028\ldots \approx 137.036$ --- the arithmetic checks; the physical derivation remains absent, exactly as the downgrade states. The $R = 2589$ normalization: the source supplies no derivation inputs in the extracted text, so its arithmetic is \textbf{flagged unverifiable-from-source} in this rebuild (carried as a downgraded claim, not re-derived).}

% ============================ PART VIII ============================
\section{PART VIII --- COMPLETE CLAIM LEDGER}
\begin{figure}[h]
\centering
\fbox{\parbox{0.85\textwidth}{\centering [figure placeholder --- original Figure 10 not recovered]\\Consolidated status across five verification cycles.}}
\caption{Figure 10. Consolidated status across five verification cycles.}
\end{figure}

\subsection{VIII.1 Established}
\begin{center}
\begin{tabular}{lll}
\toprule
claim & label & basis \\
\midrule
$D_4$ isotropic through 4th order & DERIVATION & \S II.1, symbolic + numeric \\
$D_4$ isodual; unique interference maxima & DERIVATION & \S II.2, exhaustive scan \\
Criterion $r_c = 2u^2\tau^2/(9vq_4)$ [ledger garble corrected; see front matter] & DERIVATION & \S II.3, symbolic \\
BCC recovered in 3D & DERIVATION & \S II.4, validation \\
$D_4$ selected in 4D, $S = 0.454$ & PREDICTION & \S II.5 \\
Dynamical shell selection from noise & PREDICTION & \S II.6, simulation \\
$u_{\mathrm{eff}} = u - 3v\bar{\psi}$ & DERIVATION & \S II.7, symbolic \\
$D_4$ phase band (tuned $q_0$) & PREDICTION & \S II.7 \\
Full 24-mode state is an attractor & PREDICTION & \S II.8, simulation \\
First-order transition character & DERIVATION & \S II.9 \\
$D_4 \cap \{x_4 = 0\} = \mathrm{FCC}$ & DERIVATION & \S II.10 \\
Exact dispersion, $H > 0$, CFL bound & DERIVATION & \S III.1--2 \\
Corrected envelope limit + validity & DERIVATION / PREDICTION & \S III.3 \\
$d = 4$ marginality & DERIVATION & \S III.4 \\
24-cell $= 2T \subset SU(2)$ & DERIVATION & \S IV.3, exhaustive \\
$2T/Q_8 \cong \mathbb{Z}_3$ & DERIVATION & \S IV.3 \\
McKay: $2T \leftrightarrow$ affine $E_6$ & DERIVATION & \S IV.3 \\
$E_6$ chiral & DERIVATION & \S IV.3 \\
$a_0 = 2\pi/q_0$ & DERIVATION & \S I.3 \\
$\bar{\psi}$ conserved under Model B & DERIVATION & \S I.3 \\
Point solitons from $\pi_3(S^3)$ & DERIVATION & \S V.1 \\
Independent re-derivation of load-bearing claims & CONFIRMED & \S VI-A.1 \\
Lean~4 formalization, machine-checked end to end & CONFIRMED & \S VI-A.2 \\
\bottomrule
\end{tabular}
\end{center}

\subsection{VIII.2 Excluded}
$\mathfrak{so}(8)$ as gauge algebra; SM embedding in $\mathfrak{so}(8)$; Pati--Salam from $\mathfrak{so}(8)$; $\xi_2 = -1/384$; both v92 time steps; Kerr solitons in 4D; skyrmions from a single real scalar; gauge structure internal to a single scalar; three generations as derived; monopole absence; universal optimality of $D_4$.

\subsection{VIII.3 Inputs (acknowledged assumptions)}
Generation count; the value of $\bar{\psi}$; tuning of $q_0$ onto the $D_4$ shell; conserved substrate dynamics; the dimensional assignment (\S IX.2).

% ============================ PART IX ============================
\section{PART IX --- OPEN PROBLEMS}

\subsection{IX.1 Tractable Computation}
Specified, requiring no new physics: larger-volume simulation for single-crystal ordering and phase-locking convergence; nucleation barrier height and critical droplet size; behaviour of the $D_4$ band under simultaneous detuning and higher-order terms; defect homotopy beyond elementary Burgers vectors; projection of $D_4$ to three dimensions along irrational directions.

\subsection{IX.2 Modeling Decisions}
\textbf{The dimensional assignment.} The framework carries a four-dimensional base, an internal $S^3$, and an observed 3+1 spacetime. Three readings are mutually exclusive:
\begin{itemize}
\item \textbf{(A)} the base is Euclidean spacetime, with Lorentzian signature arising dynamically. Most economical; the only reading consistent with citing four-dimensional optimality as evidence about spacetime. Cost: requires the signature mechanism of \S IX.3.
\item \textbf{(B)} the base is four spatial dimensions with one compactified. The cross-section result of \S II.10 shows this leaves FCC in three dimensions, so it is self-consistent. Cost: requires a compactification mechanism and predicts Kaluza--Klein modes; breaks the $D_4$ symmetry the selection argument uses.
\item \textbf{(C)} the base is auxiliary, with spacetime emergent from defect dynamics. Most speculative; no construction offered.
\end{itemize}
This must be decided. It is not resolvable by further calculation. The generation count, now an input, whose assumed value must be stated as such wherever fermion content is discussed.

\subsection{IX.3 Requires External Structure}
\textbf{Lorentzian signature.} Published dynamical-signature mechanisms~\cite{mukohyama2013,kehayias2014} require a one-loop effective potential with more fermionic than bosonic degrees of freedom --- the condition first identified by Greensite~\cite{greensite1993} and refined by Carlini and Greensite~\cite{carlini1994}. No purely classical Euclidean substrate is known to yield Lorentzian signature. The v92 argument from $i^2 = -1$ is invalid: canonical quantization employs complex amplitudes in Euclidean field theory as well, and an ordinary Hilbert-space inner product is positive-definite, so obtaining $g_{00} = -1$ requires a separate construction.

\textbf{Einstein field equations.} Identifying the metric with the substrate's local impedance tensor is the standard analog-gravity construction~\cite{barcelo2011} and is sound as far as it goes: wave propagation on a slowly varying background defines an effective metric, and geodesic bending and horizons follow. But the effective metric obeys the substrate's own equation of motion, not $G_{\mu\nu} = 8\pi T_{\mu\nu}$. This is a general limitation of analog models~\cite{sindoni2012}. Condition-5-type reasoning delivers effective curvature, not general relativity. Recovering the field equations, diffeomorphism invariance, the Bianchi identity, universal coupling to conserved stress-energy, and the two transverse-traceless graviton polarizations is a separate and substantially harder programme.

\textbf{Quantum mechanics beyond single-particle envelopes.} An envelope-limit Schr\"odinger equation is established (\S III.3). Reproducing entanglement and Bell~\cite{bell1964}-violating correlations from a classical substrate requires either nonlocality or superdeterminism~\cite{hossenfelder2020}. This is a boundary, not a gap.

\textbf{The monopole resolution} (\S V.3).

\subsection{IX.4 Must Be Restarted}
The $E_6$ breaking pattern and Higgs sector, which have not been derived at all; and consequently the Higgs vacuum expectation value, $A_3$ resonance mass and widths, and CKM mixing. The gauge decision makes a route available; it does not rehabilitate the prior numbers, which were computed in an excluded framework and carried independent arithmetic errors. Nothing from those calculations should be carried forward.

\subsection{IX.5 The Genuinely Open Physical Questions}
Two questions are sharper than the rest and would most advance the programme. \textbf{What fixes $\bar{\psi}$?} Under conserved dynamics it is an initial datum rather than a tuned coupling, which reclassifies the fine-tuning objection as a cosmological initial-condition question --- a different and weaker charge. But no mechanism sets its value, so whether occupying the $D_4$ band is a prediction or a tuning remains undetermined. \textbf{What supplies the nucleation event?} The transition is first order (\S II.9), so the $D_4$ phase must be nucleated. Spontaneous emergence from a featureless substrate does not occur at the relevant baseline tension.

% ============================ PART X ============================
\section{PART X --- ASSESSMENT}
Five verification cycles have produced a framework materially different from the one they began with, and smaller. Those facts are related. What has been removed was in every case removed for cause: an anisotropy coefficient not derivable from the stated lattice; a gauge algebra excluded by a Weyl-group computation; time-step values violating their own stability bound; a generation count resting on a conflation of symmetry with multiplicity; a claimed absence of monopoles inconsistent with the adopted gauge group. What remains has been computed, validated against known limiting cases, simulated where simulation was possible, and in one instance corroborated by an independent published result obtained by an unrelated method.

The reformulation from static geometry to interference dynamics was the decisive move. It removed a primitive, dissolved a standing objection at its root, supplied the selection mechanism the static formulation never had, and produced falsifiable predictions the static version could not state. That reformulation is not a reinterpretation of the same theory. It is a different and better one.

The honest position. IRH is a well-posed research programme with a verified geometric core, a defensible gauge identification, a quantified stability domain, and a specific list of open problems. It is not a completed unified theory, and the distance to one is large --- measured precisely by Part IX. Its central claim, that a high-symmetry discrete structure emerges dynamically from an oscillatory substrate rather than being imposed, is at this point supported rather than asserted. That is a smaller claim than v92 made, and it is one that has survived adversarial checking.

% ============================ APPENDICES ============================
\appendix
\section{Appendix A: Complete Record of Corrections}
Errors found across five verification cycles, including in this author's own output. Each is recorded because each altered a conclusion.
\begin{enumerate}
\item Anisotropy verification. An audit document's claimed confirmation of $\xi_2 = -1/384$ did not survive recomputation: the $D_4$ fourth moment is exactly isotropic, and the document's own substitution yields $-1/768$.
\item Real-space comparison methodology. An initial wavepacket test produced erratic 30--200\% errors, traced to a phase-sensitive metric and excessive momentum spread rather than to the derivation.
\item Robustness scan. A scan suggesting $A_4^*$ beats $D_4$ applied a near-threshold criterion outside its regime and double-counted $D_4$ against its own dual.
\item Experiment design. A simulation intended to place two shells in degeneracy placed the resonant wavenumber on a third shell. The run proved informative: that shell had zero triad content and produced lamellae, confirming the mechanism negatively.
\item Incomplete free-energy comparison. Sub-patterns within a shell were omitted; simulation exposed the omission, leading to the identification of baseline tension as control parameter.
\item Conjugacy classes of $2T$. Trace and element order merge distinct classes, giving five instead of seven; corrected by conjugation-orbit computation.
\item Weight classes in $P/Q$. An initial computation returned class 0 for all six fundamental weights, inconsistent with a nontrivial centre; corrected via Smith normal form.
\item Mixed normalization. A free-energy scan combined normalized $\tau$, $q_4$ with a raw mode count, giving 63.8\% against the correct 11.9\%.
\item Explicit integration instability. An explicit Euler step against an operator reaching $6 \times 10^4$ overflowed to NaN.
\item Conclusion contradicted by data. A script printed the pre-written conclusion ``Model B preserves the $D_4$ state'' while its own output showed the opposite. The correct explanation --- a sub-critical seed at a first-order transition --- was found only by rejecting the printed text.
\item Spectrum normalization. Power fractions computed without subtracting the mean, so nonzero $\bar{\psi}$ dominated the $k = 0$ bin.
\item Seed amplitude normalization. Seeds normalized by peak rather than per-mode amplitude, suppressing effective amplitude by a factor of 24.
\item Timestep in conserved dynamics. The explicit nonlinear term is multiplied by $k^2$ reaching 256; $dt = 0.2$ diverged.
\item Conflation of tuned and detuned $q_0$. The phase diagram used bare $r$ while simulations used detuned $q_0$; reconciling them produced the qualification of \S II.7.
\item Unstated domain convention. The simulation box length and shell normalization were never stated, making the numerics non-reproducible. Identified in independent review; fixed in \S II.6.
\item Unstated free-energy convention. $F$ as excess-over-uniform was not declared, producing an apparent 1.2\% discrepancy in independent review that turned out to be harmonic content (\S II.8).
\item Unspecified competitor shell. The key dynamical selection test in \S II.6 could not be reproduced because the triad-free competitor was never defined. Identified in independent review; now fully specified.
\item Unstated table convention. The $A_4/A_4^*$ rows are labelled by crystal while showing reciprocal-shell sizes, which reads as a transposition. Correct as computed; caption added.
\end{enumerate}
Items 1, 3, 5, 8, 10, 12 and 14 each changed a reported conclusion. Items 15--18 were reproducibility defects rather than errors of result: in each case the computation was right and the reporting was incomplete. They are recorded at the same level of prominence because a result that cannot be reproduced from the text is not, for a reader, distinguishable from one that is wrong.

\section{Appendix B: Computational Methods}
Symbolic computation used SymPy~\cite{meurer2017} throughout; the Euler--Lagrange variation, dispersion reduction, ordered-branch solution, and effective-coefficient expansion were all obtained symbolically rather than by hand. Lie-algebra results were cross-checked in Wolfram Language by generating Weyl groups from simple reflections, an algorithm structurally independent of the signed-permutation construction used in Python. Numerical optimization used SciPy~\cite{virtanen2020} (L-BFGS-B and differential evolution). Higher-order mode counting used FFT convolution on a $32^4$ grid, validated against brute-force enumeration at third and fourth order. Field simulations used a semi-implicit pseudo-spectral scheme~\cite{trefethen2000} on a $16^4$ periodic grid, with the linear operator treated implicitly in Fourier space and timesteps selected separately for conserved and non-conserved dynamics according to each scheme's stability constraint. Phase-diagram computation involved order $10^4$ independent free-energy minimizations per parameter setting. Group-theoretic results for $2T$ were obtained by exhaustive enumeration: the complete product table, conjugation orbits over all elements, and explicit coset representatives. Random seeds were fixed for reproducibility. All figures were generated directly from the computations reported, with no illustrative or placeholder data.
\flagbox{Rebuild note: the original figures were generated from computations, per this appendix, but the figure \emph{files} were not recovered from the PDF text; this rebuild's figure environments are captioned placeholders. That is a limitation of the rebuild, not of the original.}

\section{Appendix C: Terminology}
This document uses standard mathematical and physical terminology --- Lie algebra, gauge field, discrete Laplacian, dispersion relation, order parameter --- rather than framework-internal relabellings of the same objects. This is deliberate. The results above are useful only to the extent that they can be checked against the existing literature on spherical designs, Lie algebra classification, and pattern formation, and renaming standard objects would obstruct exactly that.

% ============================ REFERENCES ============================
\begin{thebibliography}{45}
\bibitem{alexander1978} S. Alexander and J. McTague, ``Should All Crystals Be bcc? Landau Theory of Solidification and Crystal Nucleation,'' \emph{Physical Review Letters} 41, 702--705 (1978). doi:10.1103/PhysRevLett.41.702
\bibitem{swift1977} J. Swift and P. C. Hohenberg, ``Hydrodynamic Fluctuations at the Convective Instability,'' \emph{Physical Review A} 15, 319--328 (1977). doi:10.1103/PhysRevA.15.319
\bibitem{mckay1980} J. McKay, ``Graphs, Singularities, and Finite Groups,'' \emph{Proceedings of Symposia in Pure Mathematics} 37, 183--186 (1980).
\bibitem{hohenberg1977} P. C. Hohenberg and B. I. Halperin, ``Theory of Dynamic Critical Phenomena,'' \emph{Reviews of Modern Physics} 49, 435--479 (1977). doi:10.1103/RevModPhys.49.435
\bibitem{delsarte1977} P. Delsarte, J. M. Goethals, and J. J. Seidel, ``Spherical Codes and Designs,'' \emph{Geometriae Dedicata} 6, 363--388 (1977). doi:10.1007/BF03187604
\bibitem{conway1999} J. H. Conway and N. J. A. Sloane, \emph{Sphere Packings, Lattices and Groups}, 3rd ed. (Springer, New York, 1999).
\bibitem{musin2008} O. R. Musin, ``The Kissing Number in Four Dimensions,'' \emph{Annals of Mathematics} 168, 1--32 (2008). doi:10.4007/annals.2008.168.1
\bibitem{korkine1877} A. Korkine and G. Zolotareff, ``Sur les formes quadratiques positives,'' \emph{Mathematische Annalen} 11, 242--292 (1877). doi:10.1007/BF01442667
\bibitem{landau1937} L. D. Landau, ``On the Theory of Phase Transitions,'' \emph{Physikalische Zeitschrift der Sowjetunion} 11, 26--47 (1937).
\bibitem{trefethen2000} L. N. Trefethen, \emph{Spectral Methods in MATLAB} (SIAM, Philadelphia, 2000).
\bibitem{virtanen2020} P. Virtanen et al., ``SciPy 1.0: Fundamental Algorithms for Scientific Computing in Python,'' \emph{Nature Methods} 17, 261--272 (2020). doi:10.1038/s41592-019-0686-2
\bibitem{bernstein1929} S. Bernstein, ``Sur les fonctions absolument monotones,'' \emph{Acta Mathematica} 52, 1--66 (1929). doi:10.1007/BF02592679
\bibitem{cohn2007} H. Cohn, J. H. Conway, N. D. Elkies, and A. Kumar, ``The $D_4$ Root System Is Not Universally Optimal,'' \emph{Experimental Mathematics} 16, 313--320 (2007). arXiv:math/0607447
\bibitem{delaat2024} D. de Laat, N. Leijenhorst, and C. de Muinck Keizer, ``Optimality and Uniqueness of the $D_4$ Root System,'' arXiv:2404.18794 (2024).
\bibitem{derrick1964} G. H. Derrick, ``Comments on Nonlinear Wave Equations as Models for Elementary Particles,'' \emph{Journal of Mathematical Physics} 5, 1252--1254 (1964). doi:10.1063/1.1704233
\bibitem{barcelo2011} C. Barcel\'o, S. Liberati, and M. Visser, ``Analogue Gravity,'' \emph{Living Reviews in Relativity} 14, 3 (2011). doi:10.12942/lrr-2011-3
\bibitem{mukohyama2013} S. Mukohyama and J.-P. Uzan, ``From Configuration to Dynamics --- Emergence of Lorentz Signature in Classical Field Theory,'' \emph{Physical Review D} 87, 065020 (2013). arXiv:1301.1361
\bibitem{kehayias2014} J. Kehayias, S. Mukohyama, and J.-P. Uzan, ``Emergent Lorentz Signature, Fermions, and the Standard Model,'' arXiv:1403.0580 (2014).
\bibitem{sindoni2012} L. Sindoni, ``Emergent Models for Gravity: An Overview of Microscopic Models,'' \emph{SIGMA} 8, 027 (2012). arXiv:1110.0686
\bibitem{borel1949} A. Borel and J. de Siebenthal, ``Les sous-groupes ferm\'es de rang maximum des groupes de Lie clos,'' \emph{Commentarii Mathematici Helvetici} 23, 200--221 (1949). doi:10.5169/seals-23242
\bibitem{humphreys1972} J. E. Humphreys, \emph{Introduction to Lie Algebras and Representation Theory} (Springer, New York, 1972).
\bibitem{pati1974} J. C. Pati and A. Salam, ``Lepton Number as the Fourth Color,'' \emph{Physical Review D} 10, 275--289 (1974). doi:10.1103/PhysRevD.10.275
\bibitem{mermin1979} N. D. Mermin, ``The Topological Theory of Defects in Ordered Media,'' \emph{Reviews of Modern Physics} 51, 591--648 (1979). doi:10.1103/RevModPhys.51.591
\bibitem{hurwitz1896} A. Hurwitz, ``\"Uber die Zahlentheorie der Quaternionen,'' \emph{Nachrichten von der Gesellschaft der Wissenschaften zu G\"ottingen}, 314--340 (1896).
\bibitem{coxeter1957} H. S. M. Coxeter and W. O. J. Moser, \emph{Generators and Relations for Discrete Groups} (Springer, Berlin, 1957).
\bibitem{gursey1976} F. G\"ursey, P. Ramond, and P. Sikivie, ``A Universal Gauge Theory Model Based on $E_6$,'' \emph{Physics Letters B} 60, 177--180 (1976). doi:10.1016/0370-2693(76)90417-6
\bibitem{fritzsch1975} H. Fritzsch and P. Minkowski, ``Unified Interactions of Leptons and Hadrons,'' \emph{Annals of Physics} 93, 193--266 (1975). doi:10.1016/0003-4916(75)90211-0
\bibitem{georgi1974} H. Georgi and S. L. Glashow, ``Unity of All Elementary-Particle Forces,'' \emph{Physical Review Letters} 32, 438--441 (1974). doi:10.1103/PhysRevLett.32.438
\bibitem{slansky1981} R. Slansky, ``Group Theory for Unified Model Building,'' \emph{Physics Reports} 79, 1--128 (1981). doi:10.1016/0370-1573(81)90092-2
\bibitem{newman1972} M. Newman, \emph{Integral Matrices} (Academic Press, New York, 1972).
\bibitem{lichnerowicz1963} A. Lichnerowicz, ``Spineurs harmoniques,'' \emph{Comptes Rendus de l'Acad\'emie des Sciences, Paris} 257, 7--9 (1963).
\bibitem{lawson1989} H. B. Lawson and M.-L. Michelsohn, \emph{Spin Geometry} (Princeton University Press, Princeton, 1989).
\bibitem{preskill1984} J. Preskill, ``Magnetic Monopoles,'' \emph{Annual Review of Nuclear and Particle Science} 34, 461--530 (1984). doi:10.1146/annurev.ns.34.120184.002333
\bibitem{demoura2015} L. de Moura et al., ``The Lean Theorem Prover (System Description),'' in \emph{Automated Deduction --- CADE-25}, Lecture Notes in Computer Science 9195, 26--48 (Springer, 2015). doi:10.1007/978-3-319-21401-6\_2
\bibitem{mathlib2020} The Mathlib Community, ``The Lean Mathematical Library,'' in \emph{Proceedings of the 9th ACM SIGPLAN International Conference on Certified Programs and Proofs (CPP 2020)}, 30--46 (2020). doi:10.1145/3372885.3373824
\bibitem{meurer2017} A. Meurer et al., ``SymPy: Symbolic Computing in Python,'' \emph{PeerJ Computer Science} 3, e103 (2017). doi:10.7717/peerj-cs.103
\bibitem{bell1964} J. S. Bell, ``On the Einstein Podolsky Rosen Paradox,'' \emph{Physics} 1, 195--200 (1964). doi:10.1103/PhysicsPhysiqueFizika.1.195
\bibitem{hossenfelder2020} S. Hossenfelder and T. Palmer, ``Rethinking Superdeterminism,'' \emph{Frontiers in Physics} 8, 139 (2020). doi:10.3389/fphy.2020.00139
\bibitem{vanmeel2009} J. A. van Meel, D. Frenkel, and P. Charbonneau, ``Geometrical Frustration: A Study of Four-Dimensional Hard Spheres,'' \emph{Physical Review E} 79, 030201(R) (2009). doi:10.1103/PhysRevE.79.030201
\bibitem{hoy2024} R. S. Hoy, ``Homogeneous crystallization in four-dimensional Lennard-Jones liquids,'' \emph{Physical Review E} 109, 044604 (2024). doi:10.1103/PhysRevE.109.044604; arXiv:2311.03465
\bibitem{greensite1993} J. Greensite, ``Dynamical origin of the Lorentzian signature of spacetime,'' \emph{Physics Letters B} 300, 34 (1993). arXiv:gr-qc/9210008
\bibitem{carlini1994} A. Carlini and J. Greensite, ``Why is spacetime Lorentzian?,'' \emph{Physical Review D} 49, 866--878 (1994). doi:10.1103/PhysRevD.49.866
\bibitem{cohn2007b} H. Cohn and A. Kumar, ``Universally optimal distribution of points on spheres,'' \emph{Journal of the American Mathematical Society} 20, 99--148 (2007). doi:10.1090/S0894-0347-06-00546-7
\bibitem{groh1999} B. Groh and B. Mulder, ``Why all crystals need not be bcc: Symmetry breaking at the liquid--solid transition revisited,'' \emph{Physical Review E} 59(5), 5613--5620 (1999). doi:10.1103/PhysRevE.59.5613
\bibitem{elitzur2011} S. Elitzur, B. Karni, and E. Rabinovici, ``A comment on BCC crystallization in higher dimensions,'' \emph{Modern Physics Letters B} 25, 1343--1347 (2011). doi:10.1142/S0217984911026358; arXiv:1101.1631
\end{thebibliography}

% ============================ REBUILD ADDENDUM ============================
\clearpage
\section*{Rebuild addendum --- verification table and flagged items (not part of v93.1)}
\addcontentsline{toc}{section}{Rebuild addendum}

\subsection*{Verification summary (this leg, 2026-11-04)}
\begin{center}
\begin{tabular}{lll}
\toprule
check set & tool + version & result \\
\midrule
Main arithmetic/algebra (59 claims) & SymPy 1.14.0 & 59 PASS, 1 intentional FAIL-flag \\
Table consistency (7 claims) & SymPy 1.14.0 & 7 PASS \\
Load-bearing second tool (20 claims) & python-flint 0.8.0 + quaternion enumeration & 20 PASS \\
Moment identities + $E_6$ det + $64/141$ & Singular 4.4.1 & PASS (differences 0; det 3) \\
Font glyph coverage & fontTools cmap & $\in$ in Fira Math; $\setminus$ fallback; {\notosans ₐ} fallback \\
\bottomrule
\end{tabular}
\end{center}
The single FAIL-flag is the source's own 11.85\% vs 11.90\% band inconsistency (\S II.7), deliberately surfaced, not a failure of the rebuild's arithmetic. Load-bearing claims verified by two tools: $S_{D_4} = 64/141$ (SymPy + flint + Singular rationals); $\tau_{D_4}/\tau_{\mathrm{BCC}} = \sqrt{2}$ (SymPy + flint squared-ratio $= 2$); ordered-branch $B_*, r_c$ (SymPy symbolic; flint confirms inputs); $D_4$ moment identities (SymPy + Singular); $2T$ closure/classes/cosets (SymPy counts + exhaustive enumeration); $E_6$ Cartan $\det = 3$ and SNF (SymPy + Singular); CFL values (SymPy + flint); $r_{\mathrm{eff}}$ values (SymPy + flint); free-energy gap (SymPy + flint).

\subsection*{Flagged unverifiable-from-source (no numbers invented)}
\begin{itemize}
\item All simulation/scan outcomes: Model A/B drift figures; resonance-scan maximum and 0.1\% gap; 400-direction spread $3\times10^{-16}$; spectral-power percentages (84\%, 100\%, 44/56\%, 77.1\%, 21.1\%, 89.5\%); phase-band percentages (11.85/11.90\%, 8.47\%); $H = 911{,}040$; attractor evolution table; phase-locking order parameter 0.26; validity-regime error table; $\beta$-window optimality scan; long-horizon stationarity figures (0.380565, 1.689598).
\item Geometric enumerations stated but not re-derived: $Q$ for $A_4^*$, $\mathbb{Z}^4$, $A_4$ (only internal consistency $Q = q_4N^2$ checked); $Q = 3384$ for $D_4$ (consistency with $q_4 = 47/8$ checked); Brillouin-zone maximum exactly 32; Weyl-group two-algorithm agreement; $w_0$ table; $E_6$ exponents and class assignments; $2T$ marks tables.
\item Lean~4 build claim (v93.1 \S VI-A.2): carried as stated; no Lean run in this leg.
\item $R = 2589$ normalization: no derivation inputs in the source text; carried as downgraded, not re-derived.
\item Original figures: not recovered; placeholders with original captions only.
\item 11.85\% vs 11.90\% band inconsistency: preserved and flagged; needs the author's clarification.
\end{itemize}

\subsection*{Still needing Brandon}
\begin{itemize}
\item Author clarification of the 11.85\% vs 11.90\% band figure (which value, or whether the difference reflects different scans).
\item The original \texttt{.tex} source and figure files, if they exist outside Drive/workspace (e.g.\ on another machine): supplying them would let a later leg replace this rebuild with a true reissue and restore real figures.
\item Confirmation that the $\tau_{D_4}/\tau_{\mathrm{BCC}} = \sqrt{2}$ correction (source says $2$) is accepted for any future v93.2; this rebuild already applies it throughout.
\item Nothing else in this leg requires Brandon: the compile, font rule, expansions, and verifications are complete as specified.
\end{itemize}

\end{document}
